% ============================================================
%  H. Høffding, Erkendelsesteori og Livsopfattelse.
%  Det Kgl. Danske Videnskabernes Selskab, Filosofiske Meddelelser II, 1.
%  København: Andr. Fred. Høst & Søn / Bianco Lunos Bogtrykkeri, 1925.
%  (Forelagt paa Mødet den 16 Oktober 1925; færdig fra Trykkeriet
%  den 7. November 1925.)
%
%  DIPLOMATIC TRANSCRIPTION (Danish)
%  Status: IN PROGRESS (begun 2026-09-30). Printed pp. 3–123: body 3–122,
%  Indhold 123. Seven chapters; 36 numbered divisions running straight
%  through the book.
%
%  Word-accuracy: to be stated here when the book is finished. The method:
%  every batch diffed against the scan's embedded layer (ocrdiff.py) AND
%  every page read at the image by a second reader who did not transcribe
%  it, with planted controls, before splicing. Until then: UNMEASURED.
% ------------------------------------------------------------
%  THE WITNESS
% ------------------------------------------------------------
%  The Royal Danish Academy's digitisation (ABBYY FineReader Server),
%  126 pp., 400 ppi colour, antiqua text layer with italics tagged.
%    Scan: ~/bibliotek/Høffding, Harald/erkendelsesteorie-livsopfattelse.pdf
%    sha256 33b0373f56088186677a8295a73467bf542d5c5f3f49c7ef414f2cf059baad4a
%
%  PAGE MAP: printed page = PDF page − 2, pp. 3–123. Verified 2026-09-30
%  from the running-head numeral of every page (PDF 6–125); p. 3 (ch. I)
%  and p. 21 (ch. II) carry no head. PDF 1–4: series half-title, series
%  notice, title leaf, blank.
% ------------------------------------------------------------

\documentclass[12pt]{article}
\usepackage[utf8]{inputenc}
\usepackage[T1]{fontenc}
\usepackage[danish]{babel}
\usepackage[protrusion=true]{microtype}
\usepackage[margin=1.4in]{geometry}
\usepackage{setspace}
\usepackage{amsmath}
\usepackage{libertinus}
\usepackage{libertinust1math}
\usepackage{textalpha}
\usepackage{fancyhdr}
\usepackage[colorlinks=true,linkcolor=black,urlcolor=black,
  pdftitle={Erkendelsesteori og Livsopfattelse},
  pdfauthor={Harald Høffding}]{hyperref}
\setlength{\headheight}{15pt}
\pagestyle{fancy}
\fancyhf{}
\fancyhead[LE,RO]{\thepage}
\fancyhead[C]{\textit{Erkendelsesteori og Livsopfattelse}}
\setlength{\parindent}{1.5em}
\setlength{\parskip}{0pt}
\setstretch{1.3}

\title{\textbf{Erkendelsesteori og Livsopfattelse}}
\author{Harald Høffding}
\date{Det Kgl.\ Danske Videnskabernes Selskab,
  \textit{Filosofiske Meddelelser} II, 1.\\
  København: Hovedkommissionær Andr.\ Fred.\ Høst \& Søn,
  Kgl.\ Hof-Boghandel.\\ Bianco Lunos Bogtrykkeri, 1925}

\usepackage{marginnote}
\usepackage{graphicx}
% The id-est sign „ɔ:“ (a turned c).
\DeclareUnicodeCharacter{0254}{\reflectbox{c}}
\newcommand{\opage}[1]{\marginnote{\footnotesize\textit{[#1]}}}
\newcommand{\apage}[1]{\marginnote{\footnotesize\textit{[#1]}}}
\begin{document}
\maketitle
\tableofcontents
\bigskip

% --- p. 3 ---
% CHAPTER HEAD as printed on p. 3: ONE centred line, "I. Erkendelsesteori."
% (with final stop), bold sans-serif; no rule; p. 3 has no running head.
% DIVISION NUMERALS (1--36, continuous) are run-in at the head of an
% indented paragraph, roman, not bold, not spaced. Unlike Totalitet, the
% print leaves NO extra space before a division (checked at "2.", p. 8,
% 300 dpi): no \medskip.
% QUOTATION MARKS: guillemets pointing inward as printed (opening mark
% points right, closing points left): opening set \guillemotright{},
% closing set \guillemotleft{}.
\section*{I.\ Erkendelsesteori.}
\addcontentsline{toc}{section}{I.\ Erkendelsesteori}

\opage{3}1. Efter den Opfattelse af Filosofi, der siden Kant's Tid
har vundet Overhaand, er Filosofi en Undersøgelse af det
menneskelige Aandslivs Forudsætninger. Menneskeligt Aandsliv
lægger sig for Dagen i Videnskab, Kunst, Moral og Religion. Paa
alle disse Omraader ligger der Forudsætninger til Grund, som ikke
behøver at fremtræde klart og med fuld Bevidsthed for dem, der
lever og arbejder paa dem. Filosofien kan hverken frembringe
Fagvidenskab, Kunst, Moral eller Religion, men den kan maaske vise,
at der paa alle disse Omraader bevidst eller ubevidst arbejdes med
Tanker, som det kan have sin Interesse at undersøge, især hvis det
kan vises, at der i disse Tanker røber sig en for menneskeligt
Aandsliv ejendommelig Struktur. Hvis det lykkes at fremdrage
saadanne Tanker, bliver Spørgsmaalet, om vi i dem har Aabenbaringer
af Tilværelsens inderste Væsen, eller om de blot staar som Former,
menneskelig Stræben paa de forskellige Omraader nu engang er henvist
til at benytte. Allerede hos Platon fremtræder den Tragik, der
indfinder sig ved ethvert Forsøg paa at paavise en absolut Grundvold
for Aandslivets Arbejde og dets Skæbner. Atter og atter forsøger
Platon at udlede al Erkendelse og al Eksistens af en højeste Ide,
der i sig selv ikke skulde behøve nogen Begrundelse. I hans Bøger
\guillemotright{}om Staten\guillemotleft{} kan man følge hans Stræben efter at naa op til, hvad der
for ham maatte være det absolute Standpunkt. Atter og atter kastes
% APPARATUS (p. 3): signature mark "1*" at the foot; not transcribed.
% --- p. 4 ---
\setcounter{footnote}{0}%
\opage{4}han tilbage i sine Tilløb for at naa dette Stade, og han
indrømmer faktisk, at han kun ved Hjælp af poetiske Analogier kan
finde Udtryk for det Maal, der stod for ham\footnote{Smlgn.
\textit{Platon's Bøger om Staten}. Analyse og Karakteristik. (Vid.
Selsk. filos. Meddelelser. 1924). p. 87---110.}. 2000 Aar senere
fastslog Kant med fuld Bevidsthed, hvad Platon saa nødig vilde
indrømme.

Skønt der foreligger betydningsfulde filosofiske Forsøg paa at
undersøge Kunstens, Moralens og Religionens Forudsætninger, kom i
Filosofien dog Undersøgelsen af Videnskabens Grundlag til at spille
den største Rolle. Erkendelsesteori er mere og mere blevet
Filosofiens centrale Fag, ja, der kan endog spores en Tendens til at
identificere Filosofi og Erkendelsesteori, saa at det bliver et stort
Spørgsmaal, hvilken Stilling de Discipliner, der drøfter de andre af
Aandslivets Omraader, kan faa indenfor Filosofien. Disse andre
Omraader stiller dog ogsaa deres Problemer, der ikke uden videre kan
behandles blot fra Erkendelsesteoriens Synspunkter. Som Tænkere,
overfor hvem et saadant Spørgsmaal maa rejses, maa især nævnes
Francis Bradley (Appearance and Reality. 1893), for hvem Filosofiens
Væsen egentlig er udtømt med Undersøgelsen af, hvilken Maalestok vi
har til at skelne mellem Realitet og Fænomen, og Hermann Cohen
(Logik der reinen Erkenntnis. 1902), Stifter af
% NOTE (p. 4): "Marburgersskolen" (double s) as printed; checked at 300 dpi.
Marburgersskolen\footnote{Smlgn. om disse Tænkere: \textit{Den nyere
Filosofis Historie}\textsuperscript{3}. IV p. 210 f. og 217 f.}.

Men der er endnu et Spørgsmaal, som i den nyeste Tid mere og mere
trænger sig frem. Af den Sindets dybe Rod, af hvilken ikke blot
Videnskab, men Kunst, Moral og Religion udspringer, fremgaar ogsaa
andre psykiske Dannelser, uvilkaarlige aandelige Krystallisationer, i
hvilke Mennesker faar Udtryk for, hvad der bevæger sig i deres Sind,
og for
% --- p. 5 ---
\opage{5}deres virkelige eller formentlige Erfaringer. Det er Livs- og
Verdensanskuelser, mere eller mindre udformede til ejendommelige
Helheder, der ikke uden videre kan henregnes under en af de i det
foregaaende nævnte fire Hovedformer for Aandsliv.

Ligesaa lidt som Filosofien kan frembringe Fagvidenskab, Kunst, Moral
og Religion, men maa se sin Opgave i at undersøge Forudsætningerne
for de Helhedsdannelser, der fremtræder paa disse Omraader, ligesaa
lidt formaar den at frembringe de individuelle Helhedsdannelser, der
fremtræder som Livs- eller Verdensanskuelser, men maa foreløbig tage
dem som noget Givet, noget Foreliggende, faktisk frembragte Synteser,
og saa bagefter søge at udfinde, ad hvilke Veje og ad hvilke Trin de
er blevne til, og hvilke Betingelser deres fortsatte Bestaaen kræver.
Den maa da opløse, hvad der uvilkaarlig er bygget op, og drøfte
dettes Betydning for menneskeligt Aandsliv som Helhed. ---

Der er mange Veje, som kan føre til Filosofi, og den enkelte Filosofs
Stilling til dens forskellige Problemer vil for en væsentlig Del
bestemmes ved den Vej, ad hvilken han først er kommen til at
filosofere. For mit Vedkommende var Interessen fra først af knyttet
til psykologiske, etiske og religiøse Spørgsmaal. Jeg bearbejdede
derfor de tre til disse Spørgsmaal svarende filosofiske Discipliner
(1876---1901). Men netop gennem Behandlingen af disse tre Problemer
førtes jeg til at se Erkendelsesproblemets fundamentale Stilling i
Forhold til dem. --- Psykologien nærmer sig dette Problem ved
Spørgsmaalet om Betingelserne for vor Tænkens Gyldighed. Baade den
gyldige og den ugyldige Tænken er jo faktisk Tænken og kan forsaavidt
undersøges psykologisk. Ved Spørgsmaalet om Forskellen mellem
Gyldighed og Ugyldighed sker Overgangen fra Psykologi til Erkendelses%
% --- p. 6 ---
\setcounter{footnote}{0}%
\opage{6}teori. Men desuden hviler selve den psykologiske Forsken, som
al anden Forsken, paa Forudsætninger, som det er Erkendelsesteoriens
Sag at drøfte. --- Ved Forsøget paa at give en Begrundelse af etiske
Domme viser det sig, at der, naar man gaar strengt logisk frem, gives
forskellige typiske Udgangspunkter, og at en erkendelsesteoretisk
Drøftelse maa sætte visse Grænser for Etikens Videnskabelighed. ---
Paa det religiøse Omraade bliver Spørgsmaalet, om Religion i nogen af
sine Former kan udvikle Tanker, der kan lægges til Grund ved
Tydningen af Tilværelsen, saaledes som denne kendes gennem Erfaring og
er tilgængelig for Videnskaben.
% APPARATUS (p. 6): a stray dot in the scan after "Videnskaben."; not transcribed.

I \guillemotright{}Den menneskelige Tanke\guillemotleft{} (1910) gaar jeg saa direkte ind paa
erkendelsesteoretisk Undersøgelse, som jeg i tre Afhandlinger (om
Totalitet, Relation og Analogi) har søgt at fortsætte, hvad enkelte
vigtige Begreber angaar.

Medens de mere konkrete Discipliner alle forudsætter Spørgsmaal, som
det bliver Erkendelsesteoriens Sag at drøfte, kan man ikke omvendt ud
fra erkendelsesteoretiske Synspunkter udlede Emnerne for de konkrete
Discipliner. Og dette gælder ikke blot for Discipliner, der staar i
nært Forhold til Filosofien, som Psykologi, Etik og
% PRINTED AS IS (p. 6): "Regionsfilosofi" (Regions-/filosofi over a line
% end) for "Religionsfilosofi"; a dropped syllable. Checked at 300 dpi.
Regionsfilosofi. Det gælder enhver speciel Videnskab.
Erkendelsesteorien forudsætter, at der er Videnskab til i Verden. Det
saa allerede Platon, naar han fandt Filosofiens Opgave i at analysere
de Forudsætninger, paa hvilke Matematiken, den eneste Fagvidenskab,
han kendte, beror\footnote{Smlgn. \textit{Platon's Bøger om Staten}.
p. 98; 107 f.}. Erkendelsesteorien har forsaavidt Karakteren af rent
logisk Analyse. Men spørges der, hvorledes det ejendommelige aandelige
Fænomen, vi kalder Videnskab, faktisk er blevet til, stilles derved en
Opgave for Psykologien paa Grundlag af Viden%
% --- p. 7 ---
\setcounter{footnote}{0}%
\opage{7}skabens Historie. Ved de Grundbegreber (Kategorier), som
videnskabelig Forsken arbejder med, viser det sig, hvorledes den
menneskelige Aands Trang og Evne udformes paa specielle Maader i
Vekselvirkning med de til hver given Tid foreliggende
Erfaringer\footnote{Smlgn. min Fremstilling af Kategorilæren i
\textit{Den menneskelige Tanke}. p. 135---245.}. Naturligvis kan og
maa selve denne psykologisk-historiske Metode igen blive Genstand for
erkendelsesteoretisk Undersøgelse. Erkendelsesproblemet opstaar stadig
paany, af den gode Grund, at ethvert Forsøg paa at behandle det,
hvilken Metode man saa end anvender, hviler paa visse Forudsætninger.

Medens psykologisk-historisk Undersøgelse indenfor Erkendelsesteorien
indtager en forberedende og underordnet Plads i Forhold til rent
logisk Analyse, forholder det sig anderledes paa de mere konkrete
Omraader, som filosofisk Tænkning kan arbejde paa. Overfor Kunst,
Moral og Religion maa der ad psykologisk-historisk Vej søges
Forstaaelse af, hvorledes saadanne Helheder, som ethvert Kunstværk,
enhver praktisk Moral, enhver levende Religion er, har kunnet blive
til, og hvorvidt de Betingelser, som herved har været afgørende, ogsaa
i Fremtiden kan antages at være tilstede. Og paa lignende Maade staar
Filosofien overfor individuelle Livsanskuelser. Ogsaa disse gør Brug
af Tanker og hviler paa visse Forudsætninger. Selv hvor en saadan
Livsanskuelse er eller mener sig at være uafhængig af al Videnskab,
vil dog Livets ubønhørlige Lov og især den indre Struktur, hvorom alt
Aandsliv vidner, kræve en vis Sammenhæng, ikke blot mellem de Tanker,
der indvæves i Livsanskuelsen, men ogsaa mellem disse Tanker og dem,
Individet ellers gør Brug af i sin Kamp for Tilværelsen. Der vil i
enhver alvorlig Livsanskuelse kunne
% --- p. 8 ---
\opage{8}spores en Trang til at finde et samlet Udtryk for de Krav, de
Opgaver og de Skæbner, Individets Liv fører med sig.

Det vil blive en særlig Opgave for denne Afhandling at undersøge
saadanne individuelle Livsanskuelsers Karakter og deres Forhold til
Spekulation (\guillemotright{}Metafysik\guillemotleft{}) paa den ene Side, Religion (især positiv,
overleveret Religion) paa den anden Side. Individuelle Livsanskuelser
er ofte, ligesom Kunst, Moral og Religion, af Filosofer og
Naturforskere blot betragtede som \guillemotright{}Antropomorfismer\guillemotleft{}, som
Aandsfænomener, der fra den strenge Tænknings Synspunkt ikke har nogen
Interesse. Allerede Parmenides skelner skarpt mellem den rene
Identitet, som Tanken søger og hævder, og \guillemotright{}de Dødeliges Meninger\guillemotleft{}, der
skyldes tilfældige og skiftende Synspunkter. Men Filosofien maa
undersøge, om der ikke bag disse Antropomorfismer kan spores et
uvilkaarligt Aandsarbejde, Vidnesbyrd om den indre Struktur, som alt
Aandsliv, selv den strengeste Videnskab, frembyder.

I min Fremstilling af den nyere Filosofi's Historie har jeg ofte
undersøgt Filosofernes individuelle Livsanskuelser i deres Forhold til
deres Systemer. I min Bog
% DEFECT (p. 8): the o of "om" carries a tail below the line (damaged
% sort or scan mark; the layer reads "Qm"); read "om". Checked at 300 dpi.
om Kierkegaard er jeg særlig kommen ind paa dette Spørgsmaal,
ligesom i min Fremstilling af Nietzsche (i \guillemotright{}Moderne Filosofer\guillemotleft{}), i
min Sammenligning mellem Kierkegaard og Nietzsche (i tredie Bind af
mine \guillemotright{}Mindre Arbejder\guillemotleft{}) og mellem Pascal og Kierkegaard (Tilskueren.
1923). Desuden har jeg i min psykologiske Studie \guillemotright{}Den store Humor\guillemotleft{}
søgt at karakterisere en ejendommelig Type for Maaden at tage Livet
paa. --- I den følgende Undersøgelse holder jeg mig til de Omraader,
jeg har tænkt mest over. Derfor kommer jeg ikke ind paa Kunstens
Problemer, saa stor personlig Betydning Kunsten ellers har haft for
mig.

2. Som jeg har søgt at vise i min Afhandling \guillemotright{}Totalitet
% --- p. 9 ---
\opage{9}som Kategori\guillemotleft{} (Vid. Selsk. Skr. 1917), spiller Helhedstanken en
dobbelt Rolle i vor Erkendelse. For det første er al Erkendelse en
Stræben efter at forbinde Iagttagelser og Forestillinger til en
Helhed. Den er en Kamp med Kaos, med alt spredt og usammenhængende.
Det gælder ikke blot om at forbinde, men om at forbinde saaledes, at
Erkendelsens Emne fremtræder i bestemte Forhold, dels mellem Emnets
egne Elementer, dels mellem Emnet som Helhed og andre Emner. Her
kundgør sig Syntese og Relation som de første i Kategoriernes Række.
Til disse første Kategorier slutter sig Kontinuitet og Diskontinuitet,
Lighed og Forskel i indbyrdes Brydning. Diskontinuitet og Forskel
vækker Tanken til Virksomhed, men Virksomheden gaar saa ud paa at
finde Kontinuitet bag Diskontinuiteten, Lighed bag Forskellene.
Saadan Brydning finder allerede Sted i den sædvanlige
Menneskeforstand, den før- og uvidenskabelige Bevidsthed. De nævnte
tre Par Begreber regner jeg derfor til de fundamentale Kategorier, der
gør sig gældende i alt Bevidsthedsliv og altsaa er fælles for den
før- og uvidenskabelige og den videnskabelige Erkendelse. De saakaldte
formale og reale Kategorier er speciellere og nøjagtigere Bestemmelser
af de fundamentale.

Muligheden af at danne fuldstændige Tankehelheder er ikke den samme
paa alle Omraader. I Logik og Matematik kan der dannes fuldkomnere
Tankehelheder end paa Omraader, hvor stadig ny Iagttagelse er
nødvendig. I en streng logisk Slutning, hvor Forholdet mellem
Forudsætninger og Slutsætning staar aldeles klart, har vi Eksempel paa
en rationel Helhed, især naar Slutningen kan vendes om, saa at man af
Slutsætningen kan udlede selve Forudsætningerne. Ethvert Led i en
saadan Tankehelhed er betinget ved de andre Led og derigennem ved selve
Helheden, og omvendt. En
% --- p. 10 ---
\setcounter{footnote}{0}%
\opage{10}saadan gennemsigtig Helhed har til alle Tider staaet som
Erkendelsens store Ideal, hvad enten man har troet, at dette Ideal
kunde virkeliggøres eller ikke. Og det er et Ideal, der hænger nøje
sammen med de første Kategorier i den fundamentale Gruppe, --- Syntese
og Relation. Den nøjeste Forbindelse mellem Tankeelementer er indgaaet
gennem de fuldt udklarede Relationer. Det var for at gøre Tilnærmelse
til dette Ideal mulig, at Platon gennem Udrensning og Potensering af de
Lighedsforhold, Iagttagelser kan frembyde, naaede til det strenge
Identitetsbegreb, han betegnede med Ordet \guillemotright{}Ide\guillemotleft{}. Og en analog
Udrensningsproces er i nyere Tid, ikke uden Indflydelse af Platon,
foregaaet med Begreberne Tal, Tid, Sted og Grad for at gøre dem
egnede til at træde i den nye Naturvidenskabs
Tjeneste\footnote{Smlg. \textit{Den menneskelige Tanke}. p.
181---201.}.

% "Streng Videnskab ..." opens an indented paragraph but fills the full
% measure; the next line ("Den sædvanlige ...") is flush left, so one
% paragraph. APPARATUS (p. 10): a scanner line runs across the page under
% it; not transcribed.
Streng Videnskab har udviklet sig af \guillemotright{}Antropomorfismer\guillemotleft{}. Den
sædvanlige Bevidsthed gaar ud fra, at Verden er, som den fremstiller
sig for os, med stor Rigdom af Kvaliteter og med en Mængde Helheder,
der staar mere eller mindre skarpt overfor hverandre. Her optræder
Helhed ikke for Videnskaben som et Ideal, der skal naas, men som et
Problem, der skal løses. Det er Videnskabens Opgave at beskrive
saadanne Helheder og derpaa at finde Betingelserne for deres Opstaaen
og fortsatte Bestaaen. Dette er den anden Side i Videnskabens Forhold
til Helhedsbegrebet.

Men Videnskaben selv udgør ikke en streng Helhed. Den deler sig i
forskellige Dele efter de forskellige Omraader, den arbejder paa, og
disse Dele optræder paa Grund af Forholdenes Natur med højst
forskellig Grad af streng Videnskabelighed. Den ene Videnskab kan
derved faa en irrationel, \guillemotright{}antropomorfistisk\guillemotleft{} Karakter i Forhold til
den anden. Man vil kunne ordne Videnskaberne i en Række,
% --- p. 11 ---
\setcounter{footnote}{0}%
\opage{11}saaledes at hvert enkelt Led staar med større Strenghed overfor
det følgende, medens det i Forhold til sin Forgænger faar en mere
ufuldkommen Karakter; --- Strenghed eller Fuldkommenhed stadig maalte
ved Tilnærmelsen til Gennemførelsen af Identitetsprincipet. Saaledes
staar Geometri i Forhold til Aritmetik, Fysik i Forhold til
Matematik, Biologi i Forhold til Fysik, Psykologi og andre
Aandsvidenskaber i Forhold til Biologi\footnote{Smlg.
\textit{Begrebet Analogi}. p. 80---88.}. Forholdet mellem det
Rationelle og det Irrationelle (Antropomorfistiske) strækker sig helt
op i Videnskabernes Række, træder ikke blot frem i Forholdet mellem
Videnskab paa den ene Side og de gængse Forestillinger paa den anden
Side. I Forholdet mellem Videnskaberne indbyrdes vil der dog kunne
gennemføres Rationalitet i Betydning af rationel Analogi. Saaledes
naar de geometriske Former opfattes som Symboler paa, hvad kun
Aritmetiken kan udtrykke med streng Nøjagtighed. Real Videnskab beror
paa Analogi mellem logisk-matematiske Rækker og Iagttagelsesrækker.
Som Kant har paavist, bestaar \guillemotright{}Anvendelsen\guillemotleft{} af logiske og matematiske
Former paa foreliggende Emner i en rationel
Analogiseren\footnote{Smlg. \textit{Begrebet Analogi}. p. 89---109.}.
Uden her igen at komme ind paa hvad jeg udførlig har drøftet i
tidligere Afhandlinger, skal jeg blot bemærke, at vort Kriterium paa
Virkelighed i Tvivlstilfælde bestaar i Muligheden af at sammenarbejde
de foreliggende Emner, saa at de viser sig at staa i rationel
Sammenhæng. Videnskabelig søger vi at gennemarbejde en saadan
Sammenhæng saaledes, at den i saa høj Grad som muligt nærmer sig til
et rent logisk eller matematisk Forhold og bliver gennemskuelig som
disse. Paa forskellige Tidspunkter i Tænkningens Historie finder vi
endog Tendenser til at gøre Virkeligheden til ren Logik eller ren
Matematik.
% --- p. 12 ---
\setcounter{footnote}{0}%
\opage{12}Émile Meyerson har nylig\footnote{\textsc{É. Meyerson}:
\textit{Déduction relativiste}. 1925. passim.} paavist en Tendens i
denne Retning hos nogle af Einstein's Tilhængere, i Lighed med, hvad
Parmenides, Descartes og Hegel forsøgte ud fra rent filosofiske
Synspunkter. ---

Trods Alt er Typen for alt Tankeliv den samme. Vi ifører os ikke i
Videnskaben en helt anden Maade at tænke paa end den, vi anvender i
vort praktiske, personlige Liv. Videnskaben er en Form for Aandsliv og
bærer alt Aandslivs karakteristiske Præg.

Den paa det af Kant givne Grundlag udviklede Erkendelsesteori, den
saakaldte Kriticisme, adskiller sig paa dette Punkt fra den saakaldte
Pragmatisme, der lægger Hovedvægten paa de ydre eller indre Motiver til
at hævde sine Tanker, uden at disse Tankers Form og Karakter skulde
staa i væsentlig Forhold til den menneskelige Aands Natur og Struktur.
Vore videnskabelige Begreber, ogsaa Grundbegreberne, skal nok svare
til menneskelig Trang, men menneskelig Evne skal ikke kunne erkendes
gennem dem. Det er jo ganske vist kun i Instinkt og Geni, at Trang og
Evne følges ad helt igennem. Men en stadig Disharmoni mellem dem vil
under Menneskeslægtens fortsatte Liv være umulig. Det er nu ikke blot
Trangen, der kan gøre sig gældende uden tilsvarende Evne; det er ogsaa
Tænkeevnen, der kan udfolde sig, uden at nogen virkelig menneskelig
Trang føles tilfredsstillet derved. I første Tilfælde forsøger Hjertets
Haab eller Frygt, dets Ønske eller dets Uvillie at diktere Tænkningen
dens Resultater. I det andet Tilfælde trænges Hjertets frie Liv tilbage
overfor virkelige eller formentlige intellektuelle Erobringer. En
Løsning af saadanne Konflikter bliver kun mulig derved, at baade
menneskelig Trang og menneskelig Evne er Udvikling underlagt. Saaledes
kan f. Eks.
% --- p. 13 ---
\setcounter{footnote}{0}%
\opage{13}Trangen til for enhver Pris at faa Løsning for alle Gaader,
fordi Livet ellers ikke skulde kunne udholdes, undergaa en Metamorfose,
der hidføres ved Erfaring om, at hvad der gør Livet værdt at leve,
ikke mindst er Tankearbejdet, Brugen af intellektuel Evne, til hvilket
Maal den saa end fører. Hvad Evnen angaar, har der ikke sjælden været
en Tendens til at tro, at man én Gang for alle skulde kunne afgøre,
hvor langt den kan føre. Endog Kant mente, at den menneskelige
Erkendelses Regnskab maatte kunne opgøres én Gang for alle. Men der er
ikke Tvivl om, at menneskelig Erkendelsesevne (\guillemotright{}Fornuften\guillemotleft{}) undergaar
Udvikling, ligesaavel som menneskelig Trang.

At alle Mennesker (Forskerne iberegnede) ikke til alle Tider har samme
\guillemotright{}Fornuft\guillemotleft{}, viser Tænkningens Historie. Et godt Eksempel haves i den
Kendsgerning, at Opdagerne af Sætningen om Energiens Bestaaen,
henimod Midten af det 19. Aarhundrede, alle hver for sig udtalte, at
denne Sætning egentlig var selvfølgelig, da der ikke er noget, som kan
blive til intet, lige saa lidt som noget kan blive til af intet. Men
spørger man, om denne Selvfølgelighed er bleven anerkendt til alle
Tider, viser Historien, at den antike Fysik var opbygget paa den
modsatte Antagelse, og at det først var det 17. Aarhundredes Tænkere,
der hævdede, at der ikke kunde være mere i Virkningen end i Aarsagen,
en Antagelse, som Gennemførelsen af Sætningen om Energiens Bestaaen
godtgør i det Enkelte.\footnote{Smlgn. \textit{Ledende Tanker i det
nittende Aarhundrede} p. 27 f.} Naar nu \guillemotright{}Selvfølgeligheden\guillemotleft{} varierer
fra Tid til Tid, vil det sige, at \guillemotright{}Fornuften\guillemotleft{} varierer. Den sunde
Menneskeforstand, af hvilken Videnskaben er udsprungen, følger
langsomt, undertiden meget langsomt, efter Videnskaben i dens Udvikling
og optager efterhaanden de mest tilgængelige videnskabelige
% --- p. 14 ---
\setcounter{footnote}{0}%
\opage{14}Resultater i sig som selvfølgelige. Der er dem, der betragter det
som selvfølgeligt, at man telegraferer gennem en Traad, men ikke
forstaar, at der ogsaa kan være traadløs Telegraferen.

Émile Meyerson giver et interessant Eksempel paa \guillemotright{}Fornuftens\guillemotleft{}
Varieren, hentet fra Udviklingen af Rumsopfattelsen. Medens man i
Oldtiden var tilbøjelig til at mene, at Retningen \guillemotright{}op og ned\guillemotleft{} indtog en
Særstilling i Forhold til de andre Retninger i Rummet, hævdede den nye
Naturvidenskab, at denne Særstilling ikke gjaldt, da Rummet var
ensartet i alle Retninger. Senere har man længe betragtet det som en
Selvfølge, at Rummet kun havde tre Dimensioner. Men der har saa udviklet
sig en ikke-euklidisk Geometri, som spiller en væsentlig Rolle ved
Fysikens nyeste Opdagelser.\footnote{\textit{Déduction Relativiste}
p. 317 f.} Meyerson mener dog, at der udenfor den geometriske Fornuft
ikke kan paavises klare Eksempler paa \guillemotright{}Fornuftens\guillemotleft{} Ændringer. Herimod
taler, mener jeg, det ovenanførte Eksempel med Sætningen om Energiens
Bestaaen. Og, idethele, hvis man opfatter Erkendelsens Grundbegreber
(Kategorierne) som Tanker, der fremdrages ved Studium af
Erkendelsesarbejdet paa dettes forskellige Trin, --- en Opfattelse, som
jeg har fremsat i \guillemotright{}Den menneskelige Tanke\guillemotleft{}, og i hvilken jeg tror, at
Meyerson er enig med mig\footnote{\textit{Déduction Relativiste} p.
170; 308 f.}, --- saa er der intet ivejen for, at Kategorier kan dø ud,
ligesom de kan blive til. De dør bort, naar Erkendelsesarbejdet overfor
nye Iagttagelser, der kræver Tydning,
% NOTE (p. 14): "nødes til slaa ind" as printed (no "at").
nødes til slaa ind paa nye Veje. En saadan hendøende, om ikke afdød
Kategori, er Substansbegrebet lige fra det Øjeblik, da Leibniz
erklærede, at vi kun kender Substanser som virkende, og at alt, hvad
der virker, er Sub%
% --- p. 15 ---
\setcounter{footnote}{0}%
% second reading 2026-09-30: the print has „Sub-/stans.“ (checked at 400 dpi); the
% transcriber had written „stanser“.
\opage{15}stans.\footnote{\textit{Den menneskelige Tanke} p. 217.} Det gælder da om at finde Loven for den i vor
Erfaring optrædende Virksomhed; gennem denne Lov har vi det
Virkendes Væsen udtrykt. Dette gælder, hvad enten der er Tale om
aandelige eller materielle Substanser. Allerede heri ligger, at det
ikke blot er paa det naturvidenskabelige, men ogsaa paa det
aandsvidenskabelige Omraade, at \guillemotright{}Fornuften\guillemotleft{} viser sig at være i
Udvikling. Ikke mindst viser det sig i den Maade, paa hvilken
nutildags religiøse og sociale Spørgsmaal behandles i alvorlig
Drøftelse. Man gaar ikke ud fra visse én Gang for alle fastslaaede
Tanker, men undersøger religiøse og sociale Fænomeners Sammenhæng med
Menneskets Natur og med de historiske Vilkaar og er klar over, at naar
disse forandres, maa ogsaa Religion og social Kultur forandres. Vor
\guillemotright{}Fornuft\guillemotleft{} er nu, i Modsætning til tidligere Tiders \guillemotright{}Fornuft\guillemotleft{}, mættet
med psykologisk og historisk Indsigt. Allerede en Sammenligning mellem
Feuerbach's og Voltaire's Drøftelser af religiøse Spørgsmaal viser
dette.

Efter denne Opfattelse af de Grundbegreber, der gør sig gældende i den
menneskelige Stræben efter Erkendelse, falder der et klart Lys over
Lessing's bekendte Sætning, at den stadige Stræben efter Sandhed er
bedre end selve Besiddelsen af Sandheden, --- at Jagten er bedre end
Byttet.\footnote{Se om Forstaaelsen af denne Sætning \textit{Den nyere
Filosofis Historie.} III\textsuperscript{3}. p. 19 f.} Lessings Begrundelse er, at kun ved Stræben, ikke ved
Besiddelse øves og udvides vore Kræfter. Og naar det nu i Erkendelsens
Historie saa ofte viser sig, at den Besiddelse, man mente at være naaet
til, var illusorisk eller i hvert Tilfælde kun havde Betydning som
Udgangspunkt for nyt Arbejde, saa træder den stadig fortsatte Stræben
% --- p. 16 ---
\opage{16}saa meget des tydeligere frem i sin Værdi. Selv om man er kommen
paa et falsk Spor, kan den konsekvente Forfølgelse af Sporet dog i
hvert Tilfælde lære, at man maa slaa ind paa andre Spor for at naa
frem. Konsekvent Tænkning bringer stedse sin Løn med sig.
Erkendelsens Ideal er, som vi har set, en rationel Helhed, vundet
gennem gensidig Bestemmelse af alle Led, og det er et Ideal, der
stedse vil stille nye Opgaver, dels af positiv, dels af negativ Art,
dels om Veje, dels om Afveje. Der er tilsidst ikke den Modsætning
mellem Stræben og Besiddelse, som Lessing gik ud fra. Den flyvende Pil
er i hvert Øjeblik baade i Hvile, idet den stedse er paa et bestemt
Sted, og i Bevægelse, da den stedse er i Begreb med at forlade dette
Sted.

3. Enhver videnskabelig Sætning svarer til et vist Udviklingstrin af
den menneskelige Fornuft, og da den sunde Menneskeforstand ikke saa
hurtig som Videnskaben kan naa dette Udviklingstrin, kan enhver
Sætning ikke uden videre gøres indlysende for den. Der er her stadig
Mulighed for Tvivl og Strid. Ofte maa for en Tid en enkelt Iagttager
eller Tænker staa som Garant for en Erkendelse, der ikke paa første
Haand kan være tilgængelig for alle. Der kan da være Tvivl, om han har
iagttaget eller tænkt rigtig. Det maa prøves --- af andre Iagttagere
eller Tænkere. Og Prøvelsen maa ske ud fra Forudsætninger, der paa det
givne Trin af Fornuftens Udvikling staar som nødvendige for
menneskelig Erkendelse. Hvis nu disse Forudsætninger bliver Genstand
for Tvivl, maa man prøve at finde andre Forudsætninger, hvis
Nødvendighed bedre kan begrundes, skønt ogsaa de naturligvis svarer
til et vist Udviklingstrin af menneskelig Fornuft. Der gives nu engang
ikke nogen absolut, men kun en hypotetisk Nødvendighed. Og det maa
% --- p. 17 ---
\setcounter{footnote}{0}%
\opage{17}herved ikke glemmes, hvad allerede Descartes gjorde opmærksom paa,
at selve det at tvivle er en Tænken, der som al Tænken hviler paa
visse Forudsætninger. Dersom disse Forudsætninger falder for Kritik,
falder ogsaa Tvivlen bort, ligesaa vel som Dogmer falder bort, naar det
Grundlag, hvorpaa de bygger, falder.

Af denne Betragtning følger, at det ikke vil være muligt én Gang for
alle at drage Grænser mellem Antropomorfismer og absolut gyldige
Forestillinger. Ogsaa den Tænken, der skulde lære os noget om \guillemotright{}Tingene
i sig selv\guillemotleft{}, er menneskelig Tænken og hviler paa visse Forudsætninger.
Jeg ser her bort fra, at ogsaa Antropomorfismer kræver at forklares.
Hvorledes er de overhovedet mulige? Spinoza har med Rette indskærpet,
at alle vore Forestillinger, de gyldige saavel som de ugyldige,
opstaar med samme Nødvendighed,\footnote{Omnes ideæ, tam adæqvatæ qvam
inadæquatæ, eadem necessitate conseqvuntur. \textit{Ethica} II, 36.
Dem.} en Sætning, der i det følgende vil finde sin Anvendelse, naar vi
føres til at se, at Filosofien ikke blot maa anvende en logisk
analyserende, men ogsaa en psykologisk-historisk Metode. Alle Opgaver
er ikke løste, fordi stedse flere Forestillinger henvises til
Antropomorfismernes Verden. Naar har man fundet den sidste
Antropomorfisme, den, som Mennesker ikke vil kunne frigøre sig for?
Den kunde kun findes, naar engang Videnskabens første Forudsætninger
og sidste Resultater kunde foreligge. Men der er her stillet en
uendelig Opgave, som Slægt efter Slægt maa arbejde paa, og som
betinger den store Horizont, under hvilken al dybere gaaende Forsken
virker.

Fra filosofisk Side er det forlængst indskærpet, at man ikke kommer ud
over Forholdet mellem Subjekt og Objekt. Selv hvis vi kunde mene at
staa overfor en absolut
% NOTE (p. 17): the note's Latin mixes "adæqvatæ", "qvam", "conseqvuntur"
% with "inadæquatæ" as printed.
% APPARATUS (p. 17): signature line "Vidensk. Selsk. Filos. Medd. II, 1." and
% signature "2" at the foot; not transcribed.
% --- p. 18 ---
\setcounter{footnote}{0}%
\opage{18}Objektivitet, maatte vi dog undersøge, hvilket Trin af
Subjektivitet, der herved forudsættes. Kant og Fichte hævdede derfor
med Rette, at den strengt videnskabelige Opfattelse af Tilværelsen kun
gælder under Forudsætning af en Subjektivitet, der har Evne til
fuldkommen Konsekvens og til at omfatte al mulig Erfaring. Intet
faktisk menneskeligt Jeg kan fyldestgøre dette Krav. Ethvert Jeg
(videnskabeligt eller uvidenskabeligt) befinder sig paa et vist Trin
af Udvikling henimod dette Ideal, der svarer til det Sandheds- eller
Virkelighedsideal, al Erkendelse er rettet imod. Det er ikke
nødvendigt at fremdrage denne Forudsætning ved ethvert enkelt
Spørgsmaal. Selv den moderne Relativitetslære drager, som rent fysisk
Teori, ikke denne Forudsætning frem, men paaviser kun, at enhver
Tidsbestemmelse, ligesaa vel som enhver Stedsbestemmelse, forudsætter
en bestemt Relation.

I sin Karakteristik af den moderne Relativitetslære lægger Émile
Meyerson stærk Vægt paa dens objektive Væsen. Naar nu alligevel
Forudsætningen om en mulig Iagttager (hvad Whitehead kalder the
percipient event) fremtræder, er det efter Meyerson's Opfattelse kun
for at slaa fast, at der her er Tale om noget, der gælder for
% LETTERSPACED (p. 18): "en-|hver" is spaced in print (300 dpi) -> \emph.
\emph{enhver} Iagttager\footnote{--- une représentation du phénomène
valable pour tous les observateurs, quels qu'ils soient.
\textit{Deduction Relativiste} p. 78.}. Et saadant rent Jeg, som Kant
og Fichte talte om, skulde ogsaa netop i hver given Situation kunne
repræsenteres af hvilket som helst Jeg, ligesom Menneskeheden til hver
given Tid og paa hvert givet Sted repræsenteres af enkelte, bestemte
Mennesker. Det rene Jeg er Menneskeslægten betragtet som en stor
Person, der stedse lærer noget.\footnote{Smlg. \textsc{Pascal}
(Fragment d'un traité du vide) Toute la suite des hommes, pendant le
cours de tant des siècles, doit être considérée comme un même homme
qui subsiste toujours et apprend continuellement.} Videnskaben er
Menneskehedens Værk, og
% PRINTED AS IS (p. 18): note 1 "Deduction" without accent (p. 12 and p. 14
% print "Déduction"); note 2 "tant des siècles" (Pascal: "tant de").
% --- p. 19 ---
\setcounter{footnote}{0}%
\opage{19}Kant's og Fichte's rene Jeg er en Personifikation af
Menneskeslægten, hvis Repræsentanter de enkelte forskende Mennesker
er. Naar Einstein konstaterer, at under visse Betingelser kan to
Iagttagere ikke blive enige om en Tidsbestemmelse, saa indtager han
selv et Standpunkt, der ligger nærmere ved \guillemotright{}det rene Jeg\guillemotleft{} end nogen af
de enkelte Iagttageres, idet han indser Grunden til, at disse ikke kan
blive enige.\footnote{Se herom \textit{Relation som Kategori} p.
74---76.}

Oswald Külpe\footnote{\textit{Die Realisirung.} Dritter Band. 1923. p.
427.} anser det for uhensigtsmæssigt at tale om et rent Jeg, dels
fordi et saadant Udtryk, hvad Romantikens Filosofi viser, let kan
tolkes metafysisk, dels fordi Begrebet \guillemotright{}rent Jeg\guillemotleft{} ikke spiller nogen
Rolle ved Begrundelsen af, at vi staar overfor en Virkelighed. ---
Hertil er at sige, at Faren for metafysisk Tydning ikke bør afholde
fra at opstille et Begreb, der viser sig at være nødvendigt, og at
Nødvendigheden i dette Tilfælde ligger i det Problem, som indeholdes,
ikke blot i Forskellene mellem samtidige Iagttagere, men ogsaa i
Forskellene mellem de skiftende Tiders Forskere. Den, der indser
Grunden til, at forskellige Forskere (samtidige eller fra forskellige
Tider) ikke er enige, indtager ligesaa vel et højere Stade, og
betegner forsaavidt en større Tilnærmelse til Menneskeslægtens
Erkendelsesideal, som den videnskabelige Virkelighedserkendelse
betegner et højere Stade end den sædvanlige Forstands Verdensbillede.

Selve den Subjektivitet, for hvilken al mulig menneskelig Viden, baade
hvad Omfang og hvad Sammenhæng angaar, kunde staa fuldt og klart,
vilde dog tænke under bestemte Forudsætninger, være et menneskeligt
Subjekt (S\textsubscript{m}). Et rent S (uden den Indeks, der betegner
visse Forudsætnin%
% APPARATUS (p. 19): signature "2*" at the foot; not transcribed. The
% second a of "saadant" (Külpe paragraph) is a damaged sort; reads "a".
% --- p. 20 ---
\setcounter{footnote}{0}%
\opage{20}ger) kan vi ikke tænke os. Men et idealiseret Menneskesubjekt
(S\textsubscript{m}) maatte besidde en Forstaaelse af, hvorfor
forskellige empiriske Subjekter (S\textsubscript{α} S\textsubscript{β}
S\textsubscript{γ}\,\dots), alle hørende ind under S\textsubscript{m},
maatte komme til forskellige Opfattelser.\footnote{Smlgn. \textit{Den
menneskelige Tanke} p. 304.} S\textsubscript{m} er selv i stadig
Udvikling, og den fulde Sandhed vil ikke blot bestaa i en objektiv
Lære, men ogsaa i en Forklaring af, hvorfor ikke alle forskende Aander
opfatter Sagerne paa den Maade, den objektive Lære angiver.

% SMALL CAPS (p. 20): "Leibniz" is set in small caps in the body here (400 dpi).
I \textsc{Leibniz}' Monadelære\footnote{Se om denne \textit{Den nyere
Filosofis Historie}\textsuperscript{3} II p. 135---149.} er dette
Synspunkt fremsat paa genial Maade, kun at han indførte den
unødvendige Antagelse, at hver Monade var en i sig afsluttet Verden, i
hvilken den hele store Verden afspejlede sig.
% GREEK (p. 20): subscripts alpha, beta, gamma read at 300 dpi; the
% print sets "S_m" etc. with roman S and a small subscript.
% APPARATUS (p. 20): a short centred rule ends ch. I below the notes;
% not transcribed.

% --- p. 21 ---
% CHAPTER HEAD as printed on p. 21: one centred line "II. Organisme."
% (with stop), bold sans-serif; p. 21 has no running head.
\section*{II.\ Organisme.}
\addcontentsline{toc}{section}{II.\ Organisme}

\opage{21}4. Naar Forholdet mellem Erkendelsesteori og Livsopfattelse
skal undersøges, vil det være lærerigt først at tage Ordet \guillemotright{}Liv\guillemotleft{} i
snævreste Betydning, altsaa om organisk Liv. Vel har vi ikke i dette
det simpleste Exempel paa Begrebet Helhed og paa de Problemer, dette
Begreb stiller. En Klode, et Bjerg, et geologisk Lag, et Atom er
ogsaa Helheder, der foreløbig kan optræde som Skranker for
Erkendelsen, men siden kan vise sig at stille mere eller mindre
betydningsfulde Opgaver for denne. Helhedsbegrebet har netop denne
dobbelt Stilling. Det standser foreløbig den logiske Analyse, der i de
formale Videnskaber er den overvejende. Navnlig stod Atomet i lange
Tider som en absolut Enhed; først i den nyeste Tid har det aabenbaret
sig som en Verden af Elektroner, og betydningsfulde Undersøgelser om,
hvad der foregaar i dets Indre, er derved bleven mulige og nødvendige.
Grunden til at Begrebet Helhed (Totalitet) bør opstilles som et
Grundbegreb, en Kategori, er netop den, at det ikke uden videre kan
føres tilbage til et Identitetsforhold, men kræver en særegen Række af
Undersøgelser. Et Mellemled mellem Identitet og Totalitet danner
Begrebet Kausalitet. Det viser sig nemlig, at hvad vi kalder Aarsag
aldrig er noget simpelt og enkelt, men er en Gruppe af positive og
negative Betingelser, og det viser sig fremdeles, at flere forskellige
Aarsagsrækker kan mødes, saa at der frem%
% --- p. 22 ---
\setcounter{footnote}{0}%
\opage{22}træder et helt System af Processer, af hvilke den ene betinger
den anden. Naar et saadant System træder op som givet Emne, staar det
erkendelsesteoretisk set, baade som Skranke (overfor rent formale
Begrebers Anvendelse) og som Opgave (der kræver Undersøgelse af
Betingelserne for dets Opstaaen og Bestaaen).\footnote{Se nærmere om
Grunden til, at Begrebet Totalitet opstilles som Kategori: \textit{Den
menneskelige Tanke.} p. 218---222. \textit{Totalitet som Kategori.} p.
18---22; 48---52.}

Der er en Analogi mellem alle Helheder, ligefra de simpleste til de
mest komplicerede. Da den Helhed, der kaldes Organisme, i mange
Henseender staar som typisk for den Rolle, Helhedsbegrebet spiller i
vor Erkendelse, holder vi os først til den. Ved Livsopfattelse mener
vi altsaa foreløbig videnskabelig Opfattelse af det organiske Liv. En
Undersøgelse af dette Emne vil være en god Forberedelse til at drøfte
Livsopfattelse i Betydning af Livsanskuelse, d. v. s. af den
ejendommelige Helhed af Tanker og Stemninger, der kan danne sig i
Menneskers Sind under Indflydelse af deres Oplevelser.

5. Fra et erkendelsesteoretisk Synspunkt maa det hævdes, at saa stærkt
et Helhedsindtryk en Organisme end kan gøre, er det dog uberettiget at
mene, at den paa Forhaand skulde være Undtagelse fra de almene
naturlige (fysiske og kemiske) Love. Hver enkelt organisk Proces, og
indenfor en saadan Proces hver Overgang fra Led til Led, maa forsøges
tydet efter de Synspunkter, Fysik og Kemi anlægger. Denne Tydning er,
idetmindste hidtil, kun lykkedes til en vis Grad, og navnlig
Forbindelsen af de mange forskellige Processer til et Helhedsliv, der
kan hævde sig som saadant, staar som et stadigt Problem. Dette Problem
løses ikke ved, som det sker i de saakaldte
% --- p. 23 ---
\opage{23}vitalistiske Teorier, at tillægge Organismen som Helhed en
særegen Kraft (en \guillemotright{}Livskraft\guillemotleft{}), der skulde holde Elementerne sammen.
Det vilde være at forklare idem per idem. Begrebet Helhed kan ikke
videnskabelig bruges til at betegne en Aarsag eller give en
Forklaring. --- I Afhandlingen \guillemotright{}Totalitet som Kategori\guillemotleft{} § 20 har jeg
kritiseret den saakaldte Vitalisme. Jeg vender her tilbage til
Problemet fra et noget andet Synspunkt.

Det er karakteristisk for den Maade, paa hvilken Drøftelsen om dette
Spørgsmaal føres, at de stridende Parter gensidig beskylder hinanden
for \guillemotright{}Metafysik\guillemotleft{}. De, der hævder, at der i det organiske Liv lægger sig
helt nye Kræfter for Dagen i Forhold til det fysisk-kemiske Omraade,
finder ofte \guillemotright{}Materialisme\guillemotleft{} hos dem, der hævder, at fysiske og kemiske
Love gælder i den organiske Natur som i den øvrige Natur. Til Gengæld
finder disse sidste ofte \guillemotright{}Spiritualisme\guillemotleft{} hos hine.
Erkendelsesteoretisk staar efter min Opfattelse Sagen saaledes, at der
ikke én Gang for alle kan sættes nogen Grænse for Anvendelsen af
fysiske og kemiske Synspunkter paa det organiske Omraade, men at paa
den anden Side saadan Anvendelse ikke faktisk kan gennemføres.
Problemet om en organisk Helheds Forhold til sine Betingelser, til de
enkelte Processer, i hvis Samvirken Helhedens Liv bestaar, kan hverken
løses ved at proklamere Mekanismen som Løsning af alle Gaader eller
ved at lade Helheden \guillemotright{}selv\guillemotleft{} gribe ind som en deus ex machina. Det er
ved denne Problemstilling, at Undersøgelse af Organismens Begreb vil
vise sig at være en Analogi til Undersøgelse af Personlighedsbegrebet
og af Livs- og Verdensanskuelser.

6. Som Forskere, der med mere eller mindre Konsekvens forsøger at
holde denne Linie, maa nævnes
% --- p. 24 ---
\setcounter{footnote}{0}%
\opage{24}J. S. Haldane, hvis Bog \guillemotright{}Mechanism, Life and Personality\guillemotleft{}
(1913) jeg allerede kunde omtale i Afhandlingen om Totalitet som
Kategori, og fra den nyeste Tid E. S. Russell (The Study of living
things. Prolegomena to a functional biology 1924). Denne sidste
Forfatter søger paa en interessant Maade at skelne mellem Biologi og
Fysiologi. Fysiologiens Metode er, hævder han, rent analytisk, opløser
Livets komplexe Fænomener i Elementer, der staar i fysisk og kemisk
Vexelvirkning med hverandre. Biologien derimod udforsker Livet i dets
umiddelbare Fremtræden. Denne Fremtræden, med den Enhed og
Koordination den udviser, maa tages som en Kendsgerning, saa meget end
dens Elementer eller Komponenter kan analyseres i det enkelte.
Biologien studerer Livet i dets frie Udfoldelse og finder netop
Trangen til fri Udfoldelse at være det Karakteristiske for Livet i
Modsætning til det Livløse. Denne Trang vil Russell betegne med det
græske Ord hormé (conatus)\footnote{Smlg. om \textit{Spinoza's} Begreb
om conatus: \textit{Spinoza's Ethica.} p. 56; 61.}. Et levende Væsen
bør først og fremmest studeres i dets naturlige Omgivelser, efter
hvilke dets Bevægelser er mere eller mindre lempede. Dette er den
funktionelle Biologis Opgave. Allerede i Vækstfænomenerne, altsaa ved
Organismens successive Tilblivelse, viser Nødvendigheden sig af at
tage Hensyn til Organismen som Helhed. Hvorledes gaar det til, at
Knokkelcellerne dannes til rette Tid og paa rette Sted? Og hvorfor er
Knokkeldannelsen overhovedet nødvendig? Spørgsmaal som disse kan kun
besvares, naar den enkelte Proces betragtes i dens Forhold til
Organismens Liv som Helhed, særlig til dens Liv i Fremtiden. Der viser
sig her en historisk Solidaritet mellem den enkelte Proces og andre
Processer indenfor
% --- p. 25 ---
\opage{25}samme Helhed. Der er altsaa her Brug for Syntese, ikke blot for
Analyse, naar Forstaaelse skal vindes.

Russell sammenligner Fysiologiens Forhold til de Erfaringer,
Friluftsbiologien kan gøre ude i Naturen, med den æstetiske Kritiks
Forhold til Kunstværket. En saadan Sammenligning ligger allerede til
Grund for Kant's \guillemotright{}Kritik der Urteilskraft\guillemotleft{}, som baade behandler
Begreberne Kunst og Geni og Begrebet Organisme. Det er de
uvilkaarlige og komplexe Fænomener i Naturen og i Aandens Verden, der
her stilles sammen. De staar som uforstaaelige, maaske paradoxale, og
det lykkes maaske aldrig den analyserende Forsken at finde en
udtømmende Forklaring af dem. Friluftsbiologi og Experimentalfysiologi
staar ofte her i skarp Modsætning til hinanden. Ude i Livet, som det
rører sig i Haver, Moser og Skove, kan Kræfterne udfolde sig i mange
forskellige Retninger, men den analyserende Fysiologi er tilbøjelig
til kun at anerkende, hvad der kan paavises experimentalt, altsaa
under visse bestemte afgrænsede Forhold. Og dog er --- hvad Russell
ikke tilstrækkelig fremhæver --- Experiment og Analyse nu engang den
Skærsild, enhver Sandhed maa gennemgaa. Sagen er den, at i den frie
Natur vil det som Regel ikke være muligt at paavise de simple og
umiddelbare Successionsforhold, da disse ikke er tilgængelige for
direkte Opfattelse. Kurt Lewin (Der Begriff der Genese 1922) har
vistnok med Rette hævdet, at det for det meste er gennem Slutning fra
Ligheder mellem forskellige Fænomener, at man naar til Antagelsen af,
at de er successive Led i en Udviklingsrække. Han forklarer derved, at
man hos frit levende Dyr saa sjælden kan slutte sig til Lovmæssighed i
Nedarvning. Paavisningen af umiddelbare Successionsrækker staar netop
som Verifikation af de Formodninger, Friluftsbiologien danner sig om
% --- p. 26 ---
\setcounter{footnote}{0}%
\opage{26}Sammenhængen mellem Livsformer. Det kommer da an paa, om man kan
danne saadanne Rækker, i hvilke ethvert Led staar som umiddelbar
Fortsættelse af et andet, saaledes som i Vækststadiernes Række hos
samme Organisme og i Generationernes Rækkefølge hos samme Gruppe
(forsaavidt denne Rækkefølge stadig kan følges). Ved Ordet
\guillemotright{}Genidentitet\guillemotleft{} (genetisk Identitet) betegner Kurt Lewin den Identitet
mellem Emner, der beror paa, at de umiddelbart afløser og fortsætter
hverandre, hører til samme genetiske Række.\footnote{Lewin's
Genidentitetsrækker svarer nærmest til, hvad jeg har kaldt identisk
varierende Forskelsrækker (Den menneskelige Tanke p. 167).}

Hvis Kurt Lewin har Ret, viser sig Nødvendigheden af
Experimentalfysiologi til Revision af de Resultater,
Friluftsbiologien (den funktionelle Biologi) mener at have
% PRINTED AS IS (p. 26): "naat" (for "naaet"), checked at 300 dpi.
naat. De to Forskningsretninger maa stadig supplere hinanden.
Helhedsopfattelse og Analyse udelukker ikke hinanden, hin skaffer
stadig Materiale til denne, og denne vil næppe nogensinde formaa at
opløse organisk Liv i en Række af rent mekaniske Led. En Analogi, som
her ligger nær (foruden den af Russell benyttede), er Forholdet mellem
Geometri og Aritmetik. Naar vi trækker en Linie, føjer vi stadig
Stykke til Stykke, men vi kan ikke godtgøre, at en Tilføjelse er
ligesaa stor som en anden. Alligevel er de geometriske Tilføjelsers
Række og Tallenes Række analoge, saaledes at der til enhver Overgang
indenfor den ene Række svarer en Overgang indenfor den anden. For
Forstaaelsen af organisk Liv vilde det være Idealet, om der til
enhver ny Synsmaade i funktionel Biologi kunde paavises en dermed
stemmende Synsmaade i experimental Fysiologi.

7. Helhedssynspunktet staar altsaa ikke i fuldstændig
% NOTE (p. 26): the note's "Den menneskelige Tanke" is roman in print
% (not italic, unlike every other citation of it); set as printed.
% --- p. 27 ---
\setcounter{footnote}{0}%
\opage{27}Modsætning til den mekaniske Opfattelse. De forholder sig som
Syntese til Analyse. Og den ene kan ikke udelukke den anden. Men
Helhedssynspunktet kan i og for sig ikke bruges til Forklaring. Ikke
blot i Biologien, hvor Vitalisme og Teleologi undertiden er fremstaaet
som definitive Gaadeløsninger, ogsaa andre Steder har man villet
benytte Helhedssynspunktet mod den paa streng Analyse beroende
mekaniske Naturopfattelse. Efter sin storartede Paavisning af
Sammenhængen indenfor vort Solsystem, --- af den Maade, paa hvilken de
forskellige Kloders Masser, Afstande, Hastigheder og Tætheder forholder
sig til hverandre, og paa hvilken Planeterne bevæger sig i Kredse om
Solen, istedetfor ifølge Tyngdeloven at falde ned i denne, --- erklærede
Newton, at hele denne vidunderlige Sammenhæng (elegantissima compages)
ikke kunde være opstaaet efter mekaniske Love.\footnote{Den nyere
Filosofis Historie\textsuperscript{3} II p. 208.} Men Analysen, og med
den Forsøgene paa at forklare Sammenføjningens Tilblivelse, kom senere,
selv om der stadig staar uløste Gaader tilbage.

Modsætningen mellem Totalitet og Mekanik, Syntese og Analyse kommer
saaledes atter og atter frem for vor Erkendelse. I de følgende Afsnit
vil vi møde dem paa Aandslivets Omraade.

8. E. S. Russell lader sig ikke nøje med den hidtil fremstillede
Betragtning. For ham bliver den Betydning, Helhedssynspunktet har, kun
tydelig, naar man som Analogi (og noget mere end blot Analogi!)
benytter et Henblik til den Helhed, det menneskelige Sjæleliv
frembyder, \guillemotright{}saaledes at der til Grund for levende
Væseners materielle Fremtræden ligger en Virksomhed af lignende Art som
(similar to) den, vi kender som vort eget Livs inderste
% --- p. 28 ---
\setcounter{footnote}{0}%
\opage{28}Realitet.\guillemotleft{} Paa dette Punkt afviger han fra
Haldane, hvem han ellers betragter som en af sine Forgængere. Vor egen
Erfaring er, mente han, den eneste Maalestok for Livsvirksomhed, vi
har, og hvad enten vi vil eller ikke, maa vi opfatte al anden
Livsvirksomhed efter dette Forbillede. Og det er da især den i vort
sjælelige Liv fremtrædende Tendens eller Stræben mod et Maal, der ikke
kan undværes, naar den funktionelle Biologi skal kunne løse sin Opgave.
--- I et Foredrag, som Russell, Aaret før hans Skrift udkom, holdt om,
hvad han kalder \guillemotright{}Psychobiology\guillemotleft{}\footnote{\textit{Proceedings
of the Aristotelian Society} 1923.}, udtalte han sig mere positivt om
sin Metode og kalder den endog ligefrem psykologisk; han udtaler
tillige her, at Problemet sandsynligvis er uløseligt, men at vi er
henviste til at gaa løs paa det fra to modsatte Ender af Skalaen, den
materielle og den psykologiske.

Det er vistnok saa, og det er allerede paavist af Kant, at vi kun
raader over to Analogier til Tydning af, hvad organisk Liv er,
Analogien med mekanisk Kausalitet og Analogien med psykisk Stræben
efter et Maal. Valget staar tilsidst mellem Mekanik og Personlighed.
Men ligesaa uberettiget som det vilde være at erklære Organismen for en
ren Mekanisme, lige saa uberettiget er en dogmatisk Anvendelse af
psykologiske Begreber ud over det Omraade, hvor Erfaring direkte nøder
dertil. Kun med stor Forsigtighed og med alt muligt Forbehold kan man
tale om ubevidst Sjæleliv. Selv ved de elementære sjælelige Ytringer
maa der være Betænkelighed ved at anvende psykologiske Forestillinger,
der gennemgaaende er hentede fra det mere udviklede
% PRINTED AS IS (p. 28): the reference mark after "Bevidsthedsliv" is
% printed "1", though the note it points to is numbered "2" at the foot
% (the page has two notes, marks 1 and 1). Checked at 300 dpi.
Bevidsthedsliv\footnote{Smlgn. \textit{Den menneskelige Tanke} p.
1---19.}. Og da nu, ifølge Russell selv, Begrebet Liv er mere
elementært end Begrebet Sjæl, vin%
% --- p. 29 ---
\opage{29}des der tilsidst intet ved at opstille Begrebet
\guillemotright{}Psykobiologi\guillemotleft{}. Det er bedre at holde sig til Udtrykket
\guillemotright{}funktionel Biologi\guillemotleft{}. Denne maa gaa ud paa at vise, hvilken
\guillemotright{}compages elegantissima\guillemotleft{} Organismen er, og saa overlade til
Fysiologien at forsøge en mekanisk Forklaring deraf. Det næste
Spørgsmaal vil da blive, om en saadan mekanisk Forklaring kan udtømme
alt, hvad funktionel Biologi kan forelægge. Man maa i Videnskaben ikke
blot gaa fra foreliggende Syntese til Analyse, men maa kunne gaa fra
denne Analyse til en ny Syntese, der gør hin første, foreliggende
Syntese forstaaelig.

Naar Russell gør Brug af psykologisk Analogi, er det ikke, fordi han
mener, at Sjælelivet optræder som et helt nyt og fremmed Princip i
Forhold til den mekaniske Natursammenhæng. Han er ikke Dualist, men
slutter sig nærmest til den spinozistisk-leibnizianske Opfattelse,
ifølge hvilken det Psykiske potentielt er tilstede selv i de mest
elementære Naturprocesser, skønt det først dukker aktuelt frem paa
% PRINTED AS IS (p. 29): "paa at vist Trin" (for "et vist"). Checked at
% 300 dpi.
at vist Trin indenfor det organiske Livs Udvikling. Det er nu ogsaa
klart, at man, hvis man overhovedet antager en Kontinuitet i
Tilværelsen, maa antage forudgaaende Muligheder for de nye
Tilværelsesformer, der opstaar. Men ogsaa her er det Analogi, der leder
os, og som ikke tør anvendes dogmatisk. Det syttende Aarhundredes
Naturfilosofi opfattede Hvile som Ligevægt mellem virkende Kræfter for
derved at komme ud over et pludseligt Spring ved Overgangen fra Hvile
til Bevægelse eller omvendt. Men at Hvile saaledes kun er potentiel
Bevægelse, kan kun godtgøres ved, at der indtræder Bevægelse, saasnart
Ligevægten mellem Kræfterne ophæves. En Mulighed, der aldrig bliver til
Virkelighed, er en Umulighed. Begrebet Potentialitet indeholder derfor
stedse et Problem,
% --- p. 30 ---
\setcounter{footnote}{0}%
\opage{30}ikke en Løsning. Den funktionelle Biologi har paa sit Omraade
den Betydning at stille Problemer, der ellers ikke ville blive stillede.
Der opstaar faktisk noget Nyt i Verden, uden at det ligestraks kan ses
at være en Omsætning i ny Form af noget, som forud var. Saaledes
opdukker Livet i Forhold til det Uorganiske (hvis det ikke, som nogle
mener, er ligesaa oprindeligt et Element i Naturen som det
Uorganiske). Og saaledes opstaar Sjælelivet i Forhold til det
organiske Liv, og indenfor Sjælelivet opstaar de mere komplekse Former
for Sjæleliv som noget nyt i Forhold til de mere elementære Former. Man
kan her ofte kun paavise et kronologisk Forhold, ikke et
Frembringelsesforhold. Med Rette talte G. H. Lewes om Emergenter som
forskellige fra Resultanter, og i sin Adresse paa British
Association's Møde 1921 optog Lloyd Morgan dette Begreb, idet han talte
om \guillemotright{}emergent evolution\guillemotleft{}\footnote{\textit{Advancement of
Science} 1921.}. Han vilde derved udtrykke, at det Nye maa tages som
saadant, ikke uden videre dogmatisk ses som uforklarligt. Uforklaret er
det foreløbig, og maaske for bestandig. Det staar som en Helhed, der
venter paa at blive analyseret. William James havde allerede anvendt
dette Begreb, naar for ham Udtrykket \guillemotright{}fri Villie\guillemotleft{} kun betød
noget Nyt, der dukkede op paa Villieslivets Omraade, uden Hensyn til,
om en Forklaring af det var mulig eller ikke.

9. Uden at ville gaa ind paa det store Spørgsmaal om Livets Opstaaen
paa Jorden har den amerikanske Kemiker Henderson\footnote{\textit{Bulletin
de la société française de philosophie.} 21. janvier 1921.} hævdet, at
Livets Opstaaen var forberedt ved, at vor Planets Udvikling meget
tidlig medførte Dannelsen af en Atmosfære, som indeholdt Vand og
Kulsyre foruden andre Stoffer, og derved frembragte de vigtigste
% --- p. 31 ---
\opage{31}Betingelser for organisk Liv. Livet er altsaa ikke indtraadt i
en Verden, der var det fremmed eller kun frembød Hindringer. Og deri
finder Henderson en Berettigelse til at betragte Tilværelsen fra et
teleologisk Synspunkt, idet de mere elementære Eksistensformer efter
deres egne Love danner et gunstigt Grundlag for højere
Eksistensformer. Udtrykket \guillemotright{}teleologisk\guillemotleft{} er Henderson ikke
selv tilfreds med, ligesaa lidt som med Udtrykket
\guillemotright{}finalité\guillemotleft{}, som han brugte, da han fremsatte sine Ideer i det
franske filosofiske Selskab. Det blev under Diskussionen her fra kemisk
Side fremhævet, at Forholdene paa Jordens Overflade i og for sig havde
kunnet muliggøre en Opstaaen af helt andre Eksistensformer end den
organiske, saa at det ikke var berettiget at tale om Forberedelse eller
Harmoni mellem elementære og højere Eksistensformer. Fra filosofisk
Side blev det indskærpet af Meyerson, at Udtrykket Teleologi eller
Finalitet kun bør bruges, naar der kan være Tale om bevidst Ordning af
Forholdene til at tjene et Maal, samt at man ikke maa appellere til
ikke-empiriske Aarsager, saalænge man ikke kan føre Bevis for, at
enhver Mulighed for at finde naturlige Aarsager er udelukket.

Trods disse berettigede Indvendinger mod den af Henderson fremsatte
Tankegang, bliver dog den Kendsgerning tilbage, at Livet, hvorledes det
saa ellers er opstaaet, har forefundet gunstige Betingelser. Ellers
vilde det naturligvis være forsvundet igen. Det faar Betydning for
Belysningen af de Emner, nærværende Afhandling beskæftiger sig med. Hvor
store Problemer Helhedsdannelsen paa alle Trin og under alle Former end
stiller vor Erkendelse, er der dog al Grund til at fortsætte
Undersøgelserne over Forholdet mellem Helhederne, deres Elementer og
deres
% --- p. 32 ---
\setcounter{footnote}{0}%
\opage{32}Betingelser. Der er maaske et irrationalt Forhold tilstede her,
ligesom mellem Kvadratets Diagonal og dets Side; men der kan, ligesom
her, findes flere og flere Decimaler. Og for vor hele Opfattelse af
Livet og Verden faar Helhedens Bestaaen og Udfoldelse en afgørende
Betydning. Til Begrebet Helhed er nemlig (som jeg tidligere har
vist)\footnote{\textit{Den menneskelige Tanke.} p. 238---240.
\textit{Totalitet som Kategori.} Kap. 5.} Begrebet Værdi nøje knyttet,
idet det betinges overalt ved Forholdet mellem Helheders Bestaaen og
deres Elementer og Betingelser.

10. Ved denne Betragtning kan der kastes Lys over et væsentligt Punkt i
Kant's Filosofi. Som antydet i det foregaaende var for ham Modsætningen
mellem Mekanisme og Organisme ikke fundamental, men skyldtes vor
Erkendelse, som her er henvist til at anvende to forskellige Metoder,
dels søge at gaa fra Delene til Helheden, dels søge at gaa fra
Helheden til Delene, uden helt at kunne gennemløbe nogen af disse to
Veje. Det var Kant's Overbevisning, at det er den samme Tingenes Orden,
der ligger til Grund for mekanisk Sammenhæng og for organiske Helheders
Opstaaen. Naturens Skønhed og Hensigtsmæssighed har vi derfor ikke Ret
til at betragte som Tilfældigheder, men vi maa i dem ligesaavel som i
den Aarsagssammenhæng, der fremtræder i Fænomenerne, se Udtryk for
Tilværelsens Natur. En saadan Opfattelse ligger allerede til Grund hos
Kant, da han i sin Ungdom bestred Newton's Ret til at erklære en
mekanisk Oprindelse af Solsystemets Harmoni for umulig, og da han
% PRINTED AS IS (p. 32): „før“ for „for“ (second reading 2026-09-30, 500 dpi).
senere i sit Hovedværk forkastede det teleologiske Bevis før Guds
Tilværelse. Og i sin \guillemotright{}Kritik der Urteilskraft\guillemotleft{} (1790) lærer han,
at Naturen kommer den menneskelige Aands Trang
% --- p. 33 ---
\setcounter{footnote}{0}%
\opage{33}og Opfattelse imøde ved at gøre Skønhed og Hensigtsmæssighed
mulige.

Ved denne Tankegang ophæves egentlig den skarpe Modsætning, Kant selv
i \guillemotright{}Kritik der reinen Vernunft\guillemotleft{} (1781) havde forudsat mellem den
menneskelige Erkendelse og \guillemotright{}Tingene i sig selv\guillemotleft{}. I de skønne og
hensigtsmæssige Naturfænomener kommer \guillemotright{}Tingene\guillemotleft{} os imøde; hvorfor
skulde de saa ikke komme os imøde, naar vi stræber efter Erkendelse?
--- Kant var ved sin Modsætning til Dogmatismen bleven ført til en
skarp Sondring mellem, hvad der ligger i den menneskelige Tænknings
Natur, og hvad der foreligger som givet og maa modtages, som det er.
Det er jo Tænkningens Selvvirksomhed, der gør Videnskab mulig. Men Kant
gik ud fra, at den menneskelige Tanke var som en Fremmed i Verden og
ikke kunde vente Støtte eller Bekræftelse fra selve Tilværelsen, som
han kaldte \guillemotright{}Ding an sich\guillemotleft{}. Dog antager han i sit Hovedværk,
at selve Tilværelsen er Aarsag ikke blot til det
\guillemotright{}Stof\guillemotleft{}, en Erkendelse har at bearbejde, men ogsaa til de
\guillemotright{}Former\guillemotleft{}, i hvilke Bearbejdelsen foretages. Men hvis det er saa,
kommer Tilværelsen os jo ogsaa her imøde! --- Den Tankegang, hvorpaa
\guillemotright{}Kritik der Urteilskraft\guillemotleft{} er bygget, er altsaa, naar man ser
nøje til, ikke væsentlig forskellig fra den, der ligger til Grund for
Kant's Erkendelsesteori, og Sandheden har, ligesom Skønhed og
Hensigtsmæssighed, sin Grund i Tilværelsens egen
Natur.\footnote{Denne Opfattelse af Kant's Filosofi har jeg allerede
fremsat i min Afhandling \textit{Kontinuiteten i Kant's filosofiske
Udviklingsgang} (Vid. Selsk. Skr. 1893). Den er kort fremstillet i
\textit{Den nyere Filosofis Historie.}\textsuperscript{3} III p. 59;
106.}
% APPARATUS (p. 33): a short rule below the note, then the foot line
% "Vidensk. Selsk. Filosof. Medd. II, 1." and signature mark "3"; not
% transcribed. Chapter II ends here.

% --- p. 34 ---
% CHAPTER HEAD as printed on p. 34 (below the running head, after a
% sinkage): ONE centred line, "III. Personlighed." (with final stop),
% bold sans-serif, as on p. 3.
\section*{III.\ Personlighed.}
\addcontentsline{toc}{section}{III.\ Personlighed}

\opage{34}11. Ved Begrebet Personlighed træder ligesom ved Begrebet
Organisme Helhedstanken i Forgrunden, og dermed optræder ogsaa de
Problemer, denne Tanke sætter i Bevægelse. Der er i Psykologiens
Historie fremtraadt Opfattelser, efter hvilke menneskeligt Aands- og
Sjæleliv forklaredes som opstaaede ved en Opsummering af de enkelte
Fornemmelser, Forestillinger, Følelser og Drifter, man kunde paavise.
Man mente, at Udviklingen paa det psykiske Omraade gik fra det Enkelte
og Simple til det Komplexe. Saadan Opfattelse modbevises allerede ved,
at det først er paa et senere Udviklingstrin, man bestemt kan skelne
mellem forskellige psykiske Elementer. Hos Barnet og den Vilde er der
allerede Koncentration, i Form af Reflex og Instinkt. Her kan ikke
skelnes mellem forskellige Sider ved Bevidsthedslivet; de virker
umiddelbart sammen. Og skal noget af de Elementer, psykologisk Analyse
mener at kunne skelne mellem, betragtes som oprindeligt, maatte det
være Villen, dette Ord taget i videste Betydning. Som et af de mest
primitive psykiske Fænomener har man, særlig fra psykopatologisk Side,
henpeget paa \guillemotright{}motoriske Helhedsfænomener\guillemotleft{} (Gadelius), hvor en
Paavirkning umiddelbart fører til Bevægelse. Reflex og Instinkt er
særegne, mere eller mindre komplexe Former for saadanne
Helhedstilstande. Bevægelsestrangen er fra først af herskende og danner
det stadige men dunkle Grundlag, ud fra hvilket, og i hvis
% --- p. 35 ---
\setcounter{footnote}{0}%
\opage{35}Tjeneste senere alt Erkendelsesliv og alt Følelsesliv kan
udvikle sig. Og i Vexelvirkning med disse kan da den oprindelige
Bevægelsestrang, der tillige er en Livstrang, en Trang til Bestaaen og
Udfoldelse, udvikles til speciellere og ejendommeligere Former. Naar
Trang er forbunden med Forestilling om et Maal, kan den kaldes Drift.
Et Sammenstød mellem forskellig Trang eller Drift kan føre til en indre
Kamp, som kan udklares gennem, hvad vi kalder Overvejelse, under
hvilken snart en, snart en anden Drift har Overvægten, indtil en enkelt
Drift med flere eller færre Eftervirkninger af Brydninger sejrer, i
løsere Form som Forsæt, i fastere Form som
Beslutning.\footnote{Smlg. min Afhandling om \textit{Begrebet Villie} i
\guillemotright{}Mindre Arbejder\guillemotleft{} III.} I hele denne Udviklingsgang, der danner en
Række, hvis Led ikke kan byttes om, er Forestillings- og
Følelseslivet væsentlig tjenende, saa betydningsfuld dens Medvirken end
kan være. --- Hvad Iagttagelse af Sjælelivet under dets Udvikling fra
Barn til Voksen eller fra det \guillemotright{}vilde\guillemotleft{} Liv til Kulturlivet saaledes
viser, bekræftes af, hvad Dyrepsykologi og Sociologi lærer os.
Thorndike har gennem sine Forsøg med Dyr paavist, at Udviklingen hos
dem gaar fra udifferentierede Tilstande til Skelnen og Generalisering,
og Sociologer som Tønnies og Levy Bruhl har paavist, at Menneskets
Sjæleliv paa tidligere Trin har en komplex Karakter, uden at vore
sædvanlige psykologiske Distinktioner dér kan anvendes.

Naar der nu ikke blot rører sig Trang eller Drift i forskellige
Retninger, men naar ogsaa Erkendelses- og Følelseslivet kan røre sig
med en vis Uafhængighed af den oprindelige Livstrang, saa bliver
Spørgsmaalet, om det sjælelige Liv kan vedblive at være en Helhed.
Helhedskarakteren beror paa den syntetiske Energi, med hvilken
Sjælelivets Elementer arbejdes sammen. Et Kaos af Elementerne
% APPARATUS (p. 35): signature mark "3*" at the foot; not transcribed.
% --- p. 36 ---
% PRINTED AS IS (p. 36): "Sindssygdom, Og" (comma before a capital, for
% a full stop). Checked at 300 dpi.
\opage{36}betegner Overgang til Sindssygdom, Og der kan derfor skelnes
mellem Syntesen som alt sjæleligt Livs Grundform, og det Indhold, der
sammenfattes med større eller mindre Energi og med større eller mindre
Held. Allerede i min Psykologi (1882) har jeg derfor skelnet mellem
% PRINTED AS IS (p. 36): "det Indhold, den sammenfattes" (den, not
% der) and "Atomisik" (Atomi-/sik, for "Atomistik"). Checked at 300 dpi.
formalt og realt Selv. Hint kundgør sig i den syntetiske Energi, dette
i den kvalitative Ejendommelighed af det Indhold, den sammenfattes.
Naar jeg saa stærkt fremhævede Syntesebegrebet, var det, fordi det
dengang gjaldt om Enheden og Sammenhængen i Bevidsthedslivet overfor
den psykologiske Atomisik, der raadede i den engelske
Associationspsykologi, som i andre Henseender har sine store
Fortjenester. Jeg henviste allerede dengang til, at ved en genetisk
Fremstilling maatte man begynde med
% PRINTED AS IS (p. 36): "Tilstande. der" (full stop for a comma).
% Checked at 300 dpi.
Tilstande. der nærmest var at betegne som Villiestilstande, naar man
tog Ordet Villen i videste Betydning. Desværre kom jeg ikke til at
omarbejde min Psykologi fra dette Synspunkt.

Særegne Problemer opstaar ved, at det sjælelige Indhold ikke altid kan
sammenfattes paa harmonisk Maade. Som allerede omtalt, kan der foregaa
Brydninger mellem forskellig Trang og mellem modstridende Drifter, saa
det gælder om at sammenarbejde dem til en Helhed, hvad der maaske kun
kan ske ved Fortrængen af enkelte genstridige Elementer. Aandelig
Udvikling begynder, ligesom den legemlige Udvikling i
Fostertilstanden, ofte ud fra spredte Punkter. Paa dette Sporadiske
lagde Sibbern stor Vægt i sin Psykologi. Det er det, der stiller os
vore Opgaver. --- Rent teoretisk stiller sig den Opgave at finde et
Enhedsbaand mellem vore forskellige Iagttagelser, og i Analogi hermed
stilles der os den Opgave at opbygge en Karakter af det Kaos af
Drifter og Tilskyndelsestrang, som mere eller mindre rører sig i
enhver af os. Hos de saakaldte hypnoide Individer
% --- p. 37 ---
\setcounter{footnote}{0}%
\opage{37}gør der sig en Disposition gældende til Kløvning af Sjælelivet,
idet der rører sig Tendenser til højst forskellige Maader at tage
Livet paa. En enkelt saadan Tendens vinder maaske Overhaand, og der kan
da ske en Helhedsdannelse omkring den som Midtpunkt. Det kan endog,
som især Flournoy\footnote{Smlgn. \textsc{Claparède}: \textit{Theodore
Flournoy.} Sa vie et son œuvre. 1921. p. 56---65.} har paavist, være
heldbringende, at der udvikler sig en saadan partiel Totalitet, da det
dermed kan blive muligt, at en vanskelig Opgave løses eller en farlig
Tilskyndelse overvindes, skønt en fuldstændig Samling af alle Sindets
% PRINTED AS IS (p. 37): "naat" (for "naaet"), "Ofter" (for "Ofte"),
% "Karakterudviking" (l dropped). Checked at 300 dpi.
Kræfter ikke er naat. Ofter dyrker Sindssygelæger hos deres Patienter
en enkelt Tendens, som har særlig Værdi, for herigennem at fortrænge
andre, mindre heldige partielle Totaliteter. Ja, man kan gaa et Skridt
videre. Ingen virkelig Karakterudviking er mulig, uden at en eller
anden Tendens eller Interesse kommer til at raade, hvilken Værdi den
saa ellers maatte have. Det er urigtigt, naar vi taler om, at vi
\guillemotright{}har\guillemotleft{} visse Egenskaber. Vore væsentlige Egenskaber \guillemotright{}har\guillemotleft{} vi
ikke, men vi \guillemotright{}er\guillemotleft{} dem. Napoleon skal have sagt, at han ikke
% NOTE (p. 37): "Ærgærrighed", "Ærgærrigheden" as printed (twice).
\guillemotright{}havde\guillemotleft{} Ærgærrighed, men at Ærgærrigheden gennemtrængte hele hans
Væsen, ligesom Blodet strømmede gennem hele hans Krop.

Hos Gustaf Fröding har den indre Kamp mellem store sjælelige
Modsætninger for en Tid ført til Vanvid. Hans Breve, især de til hans
Søster, giver herom Oplysninger af stor psykologisk Interesse. Han
taler om den forfærdelige Spaltning mellem Finfølelse og Simpelhed,
mellem Skyhed og Trods, som raader i ham. Og dertil kommer den
uhyggelige Selvkritik, der næsten gør hans kunstneriske Skaben til en
Umulighed. I sin Fortvivlelse tager han saa undertiden en Bajadsdragt
paa og udfører den til%
% --- p. 38 ---
\setcounter{footnote}{0}%
\opage{38}svarende Rolle saa godt, at han næsten bliver gal. Kom saa
hertil en absolut Tørhed og Tomhed, hvor det hverken var muligt at føle
Kærlighed eller Had, hverken Glæde eller Sorg, var det intet Under, om
Vanviddet greb ham. Og dog kunde under alt dette hans snart ophøjede,
snart sanselige Fantasier forme sig til sammenhængende Helheder.
Undertiden havde han en underlig Fornemmelse af at være en anden
Person, snart den ene, snart den anden, og det i den Grad, at ikke blot
hans Indre, men ogsaa hans Ansigtstræk og hans Skikkelse svarede til
den Person, han mente sig at være. Med Rette tilføjer han, at hvad han
saaledes har oplevet, maa være beslægtet med, hvad de saakaldte Medier
oplever. Og trods alt havde han i lysere Tider en Overbevisning om at
disse Oplevelser ikke blot betød Sygdom, men tydede paa \guillemotright{}en
Personlighedens Udviklingsproces, der længe var forberedt ved
forudgaaende Grublerier\guillemotleft{}.\footnote{\textsc{Gustav Fröding}:
\textit{Bref.} (1915). p. 206---210; 219 f.} Fuld personlig Enhed
naaede den store Digter vel ikke. Men til Gengæld har han skildret
Modsætningen mellem Lys og Mørke og mellem Drøm og Virkelighed saa
skønt og indtrængende som ingen anden Digter i vor Tid. Det er let nok
at samle sit Livsindhold i Harmoni, naar det ingen skarpe Modsætninger
frembyder. Harmoniens Værd beror paa, hvor store Disharmonier den har
haft at kæmpe med.

% PRINTED AS IS (p. 38): "sucsessivt" (for "successivt"; the same page
% has "successivt" a few lines below). Checked at 300 dpi.
Ikke blot samtidige Modsætninger kan spalte det reale Jeg og gøre
Syntesen vanskelig. Ogsaa sucsessivt optrædende Modsætninger kan gøre
saadan Virkning, navnlig medens Overgangen staar paa. Platon standsede
ved denne Vanskelighed paa Grund af den skarpe Sondring, han i den
ældre Form af sin Idelære gjorde mellem de forskellige Ideer, altsaa
mellem de forskellige Prædikater, der successivt kan tilkomme
Subjektet, eller, lad os sige, det for%
% --- p. 39 ---
\setcounter{footnote}{0}%
\opage{39}skellige Indhold, af hvilket et og samme Jeg efterhaanden kan være
optaget. Mærkeligt nok tænkte han ved dette Problem (i Dialogen
Parmenides) ikke paa, at han selv i sin geniale Lære om Eros (i
Symposion) har skildret en psykologisk Udviklingsgang, i Løbet af
hvilken Sjælen efterhaanden faar nyt Indhold.\footnote{\textit{Bemærkninger
om den platoniske Dialog Parmenides}
% PRINTED AS IS (p. 39, n. 1): the parenthesis opened at "(Vidensk." is never closed.
(Vidensk. Selsk. filos. Meddelelser (1920). p. 44---50.} Han lægger
mindre Vægt paa de psykologiske Vanskeligheder, den Smerte, der kan
fremkomme ved Overgangen fra Led til Led i denne Udvikling, end vi
vilde gøre\footnote{Denne Bemærkning skylder jeg min Hustru, der i
Samtaler ofte har kritiseret Platon's psykologiske Evolutionslære i
\guillemotright{}Symposion\guillemotleft{} fra dette Synspunkt.}, men har
som Psykolog haft en Tro paa, at en saadan Udvikling var mulig. Det er
først i nyere Tid, at Overgangenes, Krisernes Psykologi er bleven et
vigtigt Emne for Forsken, og vi kan ofte behøve Platons Tro for ikke at
tvivle om, at der kan bygges Bro over psykiske Afgrunde.

En fuldstændig Sammenfatten af Livsindholdet er ikke mulig for nogen
Personlighed. Der opdukker stedse nyt Indhold, og dermed nye Opgaver
for Sammenfatten. En absolut samlet Personlighed er et Ideal, hvad dog
ikke udelukker, at der for det opmærksomme Øje paa ethvert Trin og i
enhver Form af psykisk Liv opdages Antydninger af og Tilnærmelser til,
hvad et saadant Ideal kræver. Dette gælder baade i det uvilkaarlige og
i det fuldt bevidste Sjæleliv, skønt baade vi selv og andre bedst
forstaar, hvad der gør sig gældende i den herskende Følelse og i den
stadige Tanke. Hvad der i Øjeblikket eller hidtil i de fleste Øjeblikke
har indtaget Bevidsthedens Midtpunkt, behøver ikke at være det, der i
Virkeligheden er dybest grundet i vor Natur. Der hører megen Erfaring
og Efter%
% --- p. 40 ---
\setcounter{footnote}{0}%
\opage{40}tanke til for at blive klar over vort eget reale Jeg, og det
bliver vel egentlig kun muligt, naar Modsætning til andre Jeg'er
mærkes. Hvad det formale Jeg angaar, som bestaar i den syntetiske
Energi, der holder det sjælelige Indhold sammen, saa bliver Begrebet
derom naturligvis kun muligt overfor en forholdsvis fremskreden
intellektuel Udvikling. Fra først af kan der ikke skelnes mellem Form
og Indhold. Ordet Jeg er naturligvis ikke afgørende. Det er for Børn
foreløbig et Navn, og det kan hændes, at de kommer i Strid om, hvem af
dem der har Ret til at føre dette Navn. Filosofer er ofte bleven
beskyldte for at lægge stor Vægt paa Brugen af selve Ordet Jeg. Men
selv Fichte, der taler saa meget om \guillemotright{}das Ich\guillemotleft{},
og som ofte (saaledes nylig af en udmærket dansk Filolog) er bleven
spottet derfor, er klar over, at Tanken \guillemotright{}Jeg\guillemotleft{},
hvorved han særlig tænker paa det formale Jeg, først kan opstaa, naar
der er foregaaet en vis intellektuel Udvikling.
\guillemotright{}Saalænge jeg\guillemotleft{}, siger han\footnote{\textit{Rückerinnerungen}
§ 36.}, \guillemotright{}lever helt praktisk, helt er Liv og Handlen,
% EMPHASIS (p. 40): "min" and "mig" are letterspaced (spærret) in print.
kender jeg kun \emph{min} Følelse, Trang, Villen o. s. v. i deres
Skiften, men jeg kender ikke udtrykkelig \emph{mig} som disse
forskellige Bestemmelsers Enhed og Princip. Først naar jeg hæver mig
over disse forskellige Akters Virkelighed og med Abstraktion fra deres
Forskellighed blot sammenfatter dem som alle hørende til i mig, opstaar
Bevidstheden om Enheden som Principet for hine mangfoldige
Bestemmelser overhovedet, og dette Produkt af vor abstraherende og
sammenfattende Tænken er det, som vi kalder vor Sjæl eller
Aand.\guillemotleft{}

12. Sondringen mellem det Formale og det Reale i Personligheden er af
Betydning for Forstaaelsen af Forholdet mellem Erkendelsesteori og
Psykologi. Erkendelsesteorien
% --- p. 41 ---
\opage{41}udskiller det Formale og ser bort fra det fra Individ til
Individ og hos samme Individ i forskellige Perioder varierende reale
Indhold. Syntesebegrebet er Grundbegreb for begge Discipliner. Alle
videnskabelige Begreber bærer dets Præg. Det er intet Under, da det
videnskabelige Arbejde øves af Personligheder. Men det er tillige et
Arbejde, der fører ud over alt individuelt personligt, idet det gaar
ud paa at vinde en Forstaaelse, der kan blive fælles for alle, hvis
Tænkeevne naar en vis Udvikling (Smlgn. ovenf. § 3). En Forstaaelse
vilde være umulig, naar alle foreliggende Emner dannede en kaotisk
Forskelsrække d. v. s., en Række, hvis Led alle var lige forskellige
indbyrdes, kunde anbringes i hvilkensomhelst Orden i Forhold til
hverandre og hvert for sig tænkes borte. Det vil vel aldrig være
muligt at godtgøre, at der virkelig foreligger en saadan Række. Men den
kan bruges til at belyse Tankearbejdets Natur. Allerede den sædvanlige
Menneskeforstand søger ud over det Kaotiske, hvor det til en vis Grad
melder sig, og forbinder uvilkaarlig Emnerne paa en eller anden Maade.
Videnskaben fortsætter dette Arbejde gennem streng Prøvelse af de
enkelte Emner og af de Maader, paa hvilke de kan forbindes. Gennem en
Række af Trin, af hvilke de vigtigste i \guillemotright{}Den menneskelige
Tanke\guillemotleft{} (p. 160---173) er nærmere karakteriserede, søger
Videnskaben at reducere alle Forskelle til at være skjulte
Identiteter. Den rene Identitet staar forsaavidt som Videnskabens
Ideal. Men selve Identitetsbegrebet er, som alle Begreber, Frugt af en
Syntese, idet det betegner et Forhold mellem Emner, der hidtil stod
som forskellige. Hvor saadan Syntese lykkes, har Videnskaben vundet en
Sejr.

13. Der har været Tider, da man svælgede i Identitet som Højdepunktet
for al Videnskab. Det at føre alle Emner
% --- p. 42 ---
\setcounter{footnote}{0}%
\opage{42}tilbage til ét Emne, der stod som identisk med sig selv og
derfor ingen Forklaring behøvede, medens det tillige laa til Grund for
alle andre Emner, stod som det store Ideal, man undertiden ikke
tvivlede om at have naaet. En Grænse for Gennemførelsen af dette Ideal
fremtræder allerede i Kvalitetsforskellighederne, med deres Varieren i
Grad, Tid, Tal og Sted. Det er her de rationelle Analogier (se ovenfor
§ 2), der træder i Identitetens Sted. Men en særlig betydningsfuld
Skranke betegner her Helhedsbegrebet, som finder Anvendelse, hvor
forskellige Kvaliteter i forskellige Grads-, Tids-, Tal- og
Stedforhold fremtræder i Erfaringen som forbundne i mere eller mindre
fast Sammenhæng. En saadan Helhed er Organismen, og en saadan Helhed
har vi ogsaa i Personligheden, --- den Personlighed, i hvilken selve
Videnskaben har sit Udspring. Det mærkelige er altsaa, at der ved
personligt Arbejde er frembragt Tanker, der synes at true selve
Personlighedens Realitet. Thi naar Identitet er Udtryk for Sandhed,
kan Personlighed ikke være et sandt Begreb, dels fordi der er mange
Personligheder, dels fordi enhver enkelt Personlighed er en
ejendommelig faktisk Syntese af forskellige Emner og Elementer.
Identitet og Totalitet staar tilsyneladende her som fjendtlige Magter
overfor hinanden. Naar Totalitet alligevel er optagen som Led i
Kategoriernes Række, er det først og fremmest, fordi dette Begreb gør
opmærksom paa særegne Opgaver for Tænkningen\footnote{Se nærmere
\textit{Den menneskelige Tanke} p. 218---227 og Afhandlingen
\textit{Totalitet som Kategori}.}. Hvad Personlighed angaar, bliver
Opgaven først den at opfatte den enkelte Personlighed i dens
Ejendommelighed, leve sig ind i den, foreløbig se Tingene fra dens
Synspunkt, med andre Ord paavise den særlige og fra andre Former
forskellige Form for Tilværelse, som
% --- p. 43 ---
% APPARATUS (p. 43): a stray speck stands over the r of "her" (first word); not transcribed.
\opage{43}her træder frem. Fremdeles gælder det om at undersøge,
hvorledes alle Elementer i den personlige Helhed er paa deres Plads
dér, saa at de hver for sig bidrager til Helhedens Bestaaen og
Udvikling. Selv Disharmonier kan være paa deres Plads, hvis der uden
dem ikke kan naas en saa fuld Harmoni som muligt for Helhedens
Vedkommende. Allerede her er Analyse nødvendig. Men som tredie Opgave
fremtræder derefter den at undersøge Betingelserne for en saadan
Helheds Opstaaen og fortsatte Bestaaen. Det sker ad
psykologisk-historisk Vej, og denne Undersøgelse faar en genetisk
Karakter, idet den gaar ud paa at finde de Mellemled, gennem hvilke
Personligheden har naaet den karakteristiske Form, i hvilken den
fremtræder. Denne Opgave er i Grunden aldrig fuldt ud løselig. Det
personlige staar som en irrationel Størrelse, der kun tilnærmelsesvis
kan bestemmes ved særegne Tal. Det er i Forhold til det rene
Identitetsprincip, at Personlighed staar som noget Irrationelt (i den
angivne Betydning). Men Totalitetsbegrebet og med det særlig
Personlighedsbegrebet giver Anledning til et andet Begreb om
Irrationalitet. Ud fra dette Begreb bliver det Irrationelle det, der
ikke tilfredsstiller Grundbetingelsen for en Personligheds Bestaaen i
dens fulde Ejendommelighed og derfor vanskelig eller slet ikke kan
tænkes som vedvarende Element i den. Særlig ved Undersøgelsen af en
Personligheds Genesis kan der staa mange Gaader uløst tilbage. Men
ligesaa lidt, som naar Talen er om det organiske Liv, indeholder selve
Helheden en Forklaring. Naar man beraaber sig paa sin Personlighed, er
man ved Enden af sin Erkendelse og det samme gælder, naar man vil
forklare et andet Menneskes Personlighed ved --- hans Personlighed. Der
gives en psykologisk Vitalisme ligesom der gives en fysiologisk
Vitalisme. Der kan
% --- p. 44 ---
\setcounter{footnote}{0}%
\opage{44}opstaa Sammenstød mellem to forskellige Interesser her: den at
naa til en fyldig Karakteristik af en Personlighed i dens individuelle
Ejendommelighed og den at naa til en Forstaaelse af alle Elementer i
denne Personlighed og af den Maade, paa hvilken de har føjet sig
sammen til denne Helhed. De to Interesser strider i Grunden ikke mod
hinanden, men det er sjeldent, at de gør sig lige stærkt gældende hos
samme Forsker.

% PRINTED AS IS (pp. 44--45): "Müller-Freienfells" throughout (the author is Müller-Freienfels).
14. I sine interessante Bøger om Irrationalisme og
Individualitet\footnote{\textit{Der Irrationalismus} (1922). ---
\textit{Philosophie der Individualität} (1920).} bruger
Müller-Freienfells gennemgaaende Ordet irrationel om, hvad der ikke
gaar ind under den med Identitetsprincipet arbejdende Tænken. Fra et
saadant Standpunkt bliver Totalitetsbegrebet irrationelt. Saaledes
særlig al Individualitet. Videnskaben (Historieforskning og Psykologi)
faar da den vanskelige Opgave at gøre det, der i sig selv er
irrationelt, til Genstand for rationel Forstaaelse. --- I Modsætning
til denne Opfattelse mener jeg, at man maa opfatte Totalitet som en
selvstændig Kategori, der uden at fornægte Identitetsprincipet, som
ingen Tænkning kan undvære, begrunder en ejendommelig Maalestok, idet
den opfordrer til Undersøgelse af Forholdet mellem en foreliggende
Helhed og dens Elementer. Hermed følger, at som irrationelt kommer nu
Alt til at staa, hvad der ikke kan blive et Led af eller en Betingelse
for en vis given Helhed og dog faar fremmende eller hæmmende
Indflydelse paa denne. --- Det ses dog, at Begrebet Totalitet ogsaa
hos Müller-Freienfells ligger bag ved visse Former for Irrationalitet.
Saaledes naar det (Der Irrationalismus p. 218) erklæres, at
\guillemotright{}med Begrebet irrationel\guillemotleft{} skal der betegnes
\guillemotright{}et Positivt, en faktisk Livsmagt\guillemotleft{}, en saadan
der kan spores i Kunst, Moral og Religion. Anerkender man én Gang en
saadan Livsmagt,
% --- p. 45 ---
\setcounter{footnote}{0}%
\opage{45}saa indføres derved, mener jeg, et nyt Begreb om
Irrationalitet, der ikke blot er bestemt ved Modsætningen til
Identitetsprincipet, men ogsaa ved Modsætningen til, hvad
Helhedsprincipet i given Sammenhæng kræver. Man kan stille sig rent og
purt paa Identitetsprincipets Stade og erklære alt, hvad der ikke kan
fuldt ud begrundes ud fra det, for \guillemotright{}Dødeliges
Meninger\guillemotleft{} eller \guillemotright{}Antropomorfismer\guillemotleft{}.
Men i det Kaos, der saaledes dog bliver tilbage, fremtræder Grupper,
Helheder, som kræver Eftertanke, der ikke kan nøjes med
Identitetsprincipet.

Helhedsbegrebet er den nødvendige Forudsætning for Værdibegrebet. Om
Værdi kan kun tales ud fra en Helheds Synspunkt. Alt faar Værdi, hvad
der tjener til en vis Helheds Bestaaen og Udfoldelse. Da
Müller-Freienfells ikke har
% PRINTED AS IS (p. 45): "Helbedsbegrebet" (wrong sort b for h).
Helbedsbegrebet som Mellemled mellem Begreberne Identitet og Værdi,
forstaas det, at han nægter, at der kan gives en rationel Ordning af
Værdier. \guillemotright{}Die Werte bauen sich nicht auf Gesetze
auf\guillemotleft{} (\guillemotright{}Philosophie der
Individualität\guillemotleft{}. p. 204). Dette er i en vis Forstand
rigtigt. Man kan kun tale om Værdi i Forhold til en vis given Helhed,
men dennes Livsbetingelser vil kunne ordnes i en Række, hvis Lov
% NOTE (p. 45): "afhænger" is set with a narrow gap after "af" (kerning), not a word space.
afhænger af Karakteren af Helhedens Trang og Kaar\footnote{\textit{Den
menneskelige Tanke}. p. 346---349. --- \textit{Totalitet som Kategori}
Kap. 5.}. At her kan komme store Problemer frem, er en Sag for sig,
som vi i det næste Kapitel vil komme til at omtale. En ny Gruppe af
Irrationaliteter fremkommer jo kun, naar Totalitetsbegrebet lægges til
Grund ved Siden af Identitetsbegrebet.

15. Ligesom der arbejdes paa at forklare Organismen som en Mekanisme,
hvis enkelte Led kan forstaas ad fysisk og kemisk Vej, saaledes
forsøges det ofte at forklare Personlighed som Produkt af mere
elementære Kræfter. Det er i begge Tilfælde Helhedsbegrebet, det gaar
ud over.

% --- p. 46 ---
\opage{46}Man har ment, at hvad vi kalder Personlighed og overhovedet
alle psykiske Fænomener er blotte Virkninger, eller maaske
Sidevirkninger af de fysiologiske Processer. Om Forsøg i denne Retning
har jeg ofte haft Lejlighed til at udtale mig, senest i
\guillemotright{}Begrebet Analogi\guillemotleft{} (112---119). Efter den
Opfattelse, jeg har søgt at begrunde, er Forholdet mellem Fysiologi og
Psykologi et Eksempel paa rationel Analogi, hvorimod den ene af disse
Discipliner ikke kan begrundes ved ren Deduktion ud fra den anden.
Derved bortfalder baade Materialismen og den i det Foregaaende omtalte
\guillemotright{}Psychobiologi\guillemotleft{}. Formedelst den rationelle
Analogi mellem de to Discipliner kan fysiologisk Kundskab og Analyse
paa mange Maader dels bestyrke, dels berigtige, hvad psykologisk
Undersøgelse har ført til at antage. Man kan saaledes (som især
Bechterew har vist) i Nervesystemet og dets Funktioner skelne et
Centralt og et Periferisk. Der kan paavises Centralorganer, i hvilke
Eftervirkninger af tidligere Paavirkninger stadig kan ligesom
opbevares i et Slags \guillemotright{}Samlebækken\guillemotleft{} for de
Erfaringer, der er vigtigst for Organismen, saa at øjeblikkelige
Paavirkninger ikke bliver aldeles afgørende for Organismens
Reaktioner. Men naar Bechterew kalder disse Centralorganer for
\guillemotright{}den personlige Sfære\guillemotleft{}, er det tydeligt, at
han bygger paa en fra Psykologien hentet Analogi. Der kan gennem
psykologisk Iagttagelse meget vel, uden Benyttelse af nogen
fysiologisk Indsigt, paavises en Forskel mellem et Centralt og et
Periferisk. Det Centrale, som jeg plejer at kalde det reale Jeg,
betragter vi som mere karakteristisk for Personligheden end de i
enkelte Øjeblikke opdukkende Fornemmelser og Tilskyndelser. --- Hvor
langt den rationelle Analogi mellem Fysiologi og Psykologi kan
gennemføres, kan kun fortsat Erfaring afgøre. Foreløbig er der ikke
Grund til at tro andet end
% --- p. 47 ---
\setcounter{footnote}{0}%
\opage{47}at alle psykiske Tilstande og Funktioner er knyttet til
fysiologiske Tilstande og Funktioner, selv om man ikke altid kan
paavise disse. Paa Forhaand kan der ikke sættes nogen Grænse for, hvad
vi herigennem kan lære om vort Sjælelivs Vilkaar.

Et andet Forsøg paa at godtgøre Personlighedsbegrebets sekundære
Karakter blev gjort indenfor selve Psykologien, idet den ældre
engelske Skole i Forestillingsassociationens Mekanik saa den Vej, ad
hvilken Personligheden som Helhed blev til gennem Vekselsvirkning
mellem de enkelte psykologiske Elementer. Men det viser sig, at alle
Former for
% PRINTED AS IS (p. 47): "Forestillingassociationen" (no linking s) and "Forstillingsindholdet" (no e).
Forestillingassociationen kan føres tilbage til en Tendens til at samle
Forstillingsindholdet til en Helhed, og denne Helhedsassociation
vidner om, at Trang og Evne til Sammenfatten stedse ytrer sig. Det er
Psykologiens Opgave at paavise en syntetisk Energi som virksom i alle
Sjælelivets Fænomener fra de mest elementære til de mest komplexe.
Hvad vi kalder Personlighed, er kun en udpræget Form for og Resultat
af saadan psykisk Energi\footnote{Foruden at henvise til min
\textit{Psykologi} kan jeg her henvise til den Undersøgelse af
\guillemotright{}Totalfølelsers\guillemotleft{} Opstaaen, jeg har
anstillet i min Bog \textit{Den store Humor}.}.

Hverken hvad Personligheden eller hvad Organismen angaar, kan Livets
Syntese fuldstændig forklares ad Analysens Vej. Der vil her stedse
staa en Kamp mellem to forskellige Principer, det ene, som gaar ud fra
Elementerne og vil forklare Helheden som blot Sammensætning af dem,
det andet, som søger at finde Helhedspræget overalt, selv i de
simpleste organiske og psykiske Fænomener.
% APPARATUS (p. 47): a short centred rule closes the chapter below the footnote; not transcribed.

% --- p. 48 ---
\section*{IV.\ Livsanskuelse.}
\addcontentsline{toc}{section}{IV.\ Livsanskuelse}

\opage{48}16. Aandsvidenskab bliver kun mulig ved, at det personlige Liv
udfolder sig frit og naturlig. Kun derved faar Aandsvidenskaben Emner
at drøfte. Det gælder særlig for de filosofiske Discipliner
(Erkendelsesteori, Psykologi, Etik og Religionsfilosofi). Der maa
komme Forsøg paa Videnskab, for at Erkendelsesteorien kan faa et Emne
at undersøge. Sjælelivet maa udfolde sig uvilkaarlig, for at
psykologisk Iagttagelse og Analyse skal blive mulige. Og der maa
udvikle sig et etisk og religiøst Liv, for at etiske og religiøse
Problemer kan stilles og drøftes. Og saadanne uvilkaarlige Udfoldelser
vil, som Erfaring viser, føre til Dannelse af aandelige Helheder. Som
allerede omtalt ligger der en Helhedslov til Grund ved al
Forestillingsassociation. De fleste Genkendelsesakter bestaar, især
naar det gælder hurtig Besindelse, i at en Del (f. Eks. en enkelt
Egenskab) umiddelbart genkendes og Resten af Emnets Egenskaber
uvilkaarlig forudsættes. Ved saadan partiel Genkendelse bliver
Sanseillusion let mulig. Der foretages i saa Tilfælde en uvilkaarlig
Syntese, som kan berigtiges ved en
% PRINTED AS IS (p. 48): "senerc" (c for e, checked at 400 dpi).
senerc Analyse. Det er et velbekendt psykologisk Fænomen, der kan
tjene som et simpelt Eksempel paa Forholdet mellem uvilkaarlig Syntese
og en følgende Analyse. Som Levy Bruhl har paavist, er der
% UNSURE (p. 48): accent on "én" is a worn dot-like mark; read as acute (rejected: plain "en").
én primitiv Tendens til at gøre Del til Helhed, det vil sige, at
betragte et enkelt Elements Tilstedeværelse som Symptom paa, at et
Kompleks af for%
% --- p. 49 ---
\setcounter{footnote}{0}%
\opage{49}skellige Elementer foreligger. Et karakteristisk Eksempel paa,
hvad Levy Bruhl kalder Participation, er, at de Vilde tror, at man ved
at have et Billede af en Ting har hele Tingen i sin
Magt\footnote{I \textit{Begrebet Analogi} p. 7---23 har jeg nærmere
omtalt dette og tillige søgt at vise, at det her er et psykologisk
Fænomen, der kan give Forklaring af et sociologisk Fænomen, ikke
omvendt.}.

I Fænomener som de nu antydede har vi et Udgangspunkt for Forstaaelse
af Livsanskuelsers Tilblivelse, maaske, som senere vil blive omtalt,
ogsaa for Forstaaelse af Metafysik og Religion. Foreløbig holder vi os
til de individuelle Maader at opfatte og tage Livet paa, altsaa til
Synteser, der i en vis Uafhængighed af Videnskab og Religion
uvilkaarlig opstaar hos enkelte Individer. Det var en Ejendommelighed
for det 19. Aarhundrede, at det af en Række fremragende Tænkere blev
hævdet, at det ved al Tro og Livsanskuelse først og fremmest kommer an
paa Ægtheden, d. v. s. paa, at Troen virkelig var fremgaaet af
personlig Trang under Indflydelse af de Livserfaringer, der var faldet
i Individets Lod. Carlyle aabner Rækken med sin \guillemotright{}Sartor
Resartus\guillemotleft{} (1836); her i Norden optræder kort efter (men
uden Indflydelse af Carlyle) Erik Gustaf Geijer, og derpaa danske
Tænkere som Poul Møller, Sibbern og
% PRINTED AS IS (p. 49): the reference mark after "Kierkegaard" is printed "1", though it keys note 2.
Kierkegaard\footnote{Nærmere herom i \textit{Ledende Tanker i det
nittende Aarhundrede} p. 137---142.}. I Tyskland har Nietzsche's
blivende Betydning vist sig at bero paa hans Fordring om Ægthed og
Selvstændighed i Livsanskuelse. Ogsaa Dilthey har, paa Grundlag af
sine indgaaende Studier over filosofiske Grænsespørgsmaal, virket i
samme Retning, og i den nyere Tid er især Georg Simmel at nævne, som i
det mesterlige lille Skrift \guillemotright{}Der Konflikt der modernen
Kultur\guillemotleft{} (1921) med stor Klarhed har belyst Tidens Uro.
% APPARATUS (p. 49): signature line "Vidensk. Selsk. Filosof. Medd. II, 1." and "4" at the foot; not transcribed.
Medens i
% --- p. 50 ---
% NOTE (p. 50): "(élan de vie)" as printed.
\opage{50}Frankrig Bergson fra sin aandfulde Paavisning af Livets
Udfoldelsestrang (élan de vie) straks springer over i Metafysik, blev
Simmel staaende paa Tærskelen. Selv om det skulde være hans tidlige
Død, der hindrede ham i at forsøge en Vandring paa den anden Side af
Grænsen, saa er det dog karakteristisk for ham, at han betænkte sig saa
længe. Det var hans Studier over Schopenhauer og Nietzsche, der førte
ham til at lægge saa stor Vægt paa Livsanskuelsens Udspring af det
personlige Livs inderste Kilde. I sit ovennævnte Skrift viser han paa
% PRINTED AS IS (p. 50): "i hvilket" (for "i hvilke").
en interessant Maade, hvorledes saa mange af de Former, i hvilket
Livet hidtil havde lagt sig til Rette, er i Opløsning, og hvorledes
Livet i sin Utaalmodighed over, at ingen nye Former frembyder sig, hos
mange Mennesker rører sig med Fjendskab
% PRINTED AS IS (p. 50): "over-/or" across the line break, i.e. "overor" (f dropped).
overor al Form. Det er vistnok karakteristisk for det aandelige Liv i
Tyskland for Tiden, at man higer efter en Livsanskuelse, der først og
fremmest bærer Personlighedens eget Præg, saaledes at Videnskab og
Religion, om de virker med, staar som underordnede Elementer. I sin
Undersøgelse af de forskellige Livsanskuelser lægger man derfor særlig
Vægt paa det psykologiske Grundlag, der tilsidst bærer den. Udførlig
har Karl Jaspers behandlet dette Emne (Psychologie der
Weltanschauungen. 2. Ausg. 1922). Han erklærer Kierkegaard og
Nietzsche for de største
\guillemotright{}Verdensanskuelsespsykologer\guillemotleft{}, fordi de i dyb
og ejendommelig Erfaring har oplevet Tilværelsens Problemer, og fordi
de i første Linie betoner Ægtheden af det personlige Liv, i hvilken
Retning det saa ellers bevæger sig. Og de to nævnte Tænkere bevæger
sig jo ogsaa i meget modsatte Retninger. --- Man beskæftiger sig
overhovedet i Tyskland meget med Kierkegaard, efterat der nu er kommen
en fuldstændig tysk Oversættelse af hans
% --- p. 51 ---
\setcounter{footnote}{0}%
\opage{51}Værker\footnote{Spanieren \textit{Unamuno} er i sin mærkelige Bog \textit{The
tragic sense of life} (engelsk Oversættelse 1921) ogsaa paavirket af
Kierkegaard, hvem han kalder sin \guillemotright{}aandelige Broder\guillemotleft{}. Han læser ham paa
Dansk.}. Og det er mere hans Paastand, at Subjektiviteten er
Sandheden, der interesserer, end hans Polemik mod officiel Kristendom.

Jaspers' Sammenstilling af Kierkegaard og Nietzsche giver mig
Anledning til at omtale et Foredrag, jeg i 1911 holdt i Helsingfors,
og som senere tryktes i tredie Bind af mine \guillemotright{}Mindre Arbejder\guillemotleft{}. Det
giver en Sammenligning mellem Kierkegaard og Nietzsche. Jeg gør først
opmærksom paa den store Modsætning mellem de to Tænkere, men
fremhæver, at Floder, der strømmer i forskellig Retning, godt kan have
deres Udspring paa det samme Bjerg, og at Filosofien nu engang mere
interesserer sig for Kilderne end for Løbet, mere for de oprindelige
Livserfaringer end for deres Udtydning. Fælles for Kierkegaard og
Nietzsche er, at de begge ser langt ud over Livets sædvanlige,
afgrænsede, rationelle Baner, --- at de begge hævder det Inderste og
Oprindelige i Mennesket som noget, der ikke kan gøres rationelt,
medens det dog kan besidde den højeste Kraft og Værdi. De er
Symptomer paa en Trang, der rører sig i Tiden, og de giver den et
vulkansk Udtryk. Jeg søger derefter at paavise en hel Række af
Analogier mellem de to saa forskellige Livstænkere, men kan her ikke
gaa nærmere ind paa dem. Afhandlingen slutter med de Ord, som jeg
stadig kan fastholde: Saa længe Livet skrider frem, vil dets hele
Indhold ikke helt og holdent kunne lægges til Rette af Tanken, men det
vil stedse gælde om at holde Sindet aabent for det Nye, som Livet og
Erfaringen kan bringe. --- Senere har jeg anstillet en lignende
Undersøgelse af Forholdet mellem
% APPARATUS (p. 51): signature mark "4*" at the foot; not transcribed.
% --- p. 52 ---
\opage{52}Pascal og Kierkegaard (i \guillemotright{}Tilskueren\guillemotleft{} 1923). Her var
Analogierne snevrere, men Forskellene, der dels skyldtes Modsætningen
mellem Katolicisme og Protestantisme, dels de forskelligartede
Tidsforhold ved de to store Aanders Optræden, er derfor ikke mindre
karakteristisk.

Hvor megen Interesse Filosofien end kan og maa have for individuelle
Forsøg paa at tage Stilling til Livsproblemerne, saa hører det dog
netop til dens Opgaver at vogte den rationelle Tænknings Grænse, for
at Dogmatismen ikke skal holde sit Indtog under Navn af Livsanskuelse
eller Livsfilosofi. Men først og fremmest maa den søge, med de Midler,
der staar til dens Raadighed, at blive klar over de ejendommelige
psykiske Helhedsdannelser, der kan træde frem i Form af
Livsanskuelser.

17. Der er to forskellige Aarsager til, at Opmærksomheden mere og mere
har rettet sig mod individuelle Livsanskuelser.

Den ene Aarsag var man allerede i Oldtiden klar over. Den ligger i al
Videnskabs almene Karakter, som Aristoteles saa stærkt fremhævede.
Videnskaben søger Forstaaelse gennem Tilværelsens almindelige Love. Den
søger Udgangspunkter, der gør det muligt at drage Slutninger, som
gælder for alt og alle. Men Livet, og hele Tilværelsen, aabenbarer sig
for os ud fra spredte, indbyrdes meget forskellige Udgangspunkter,
hvert betinget ved sine ejendommelige Forhold. De enkelte Ting eller
Væsener fremtræder hver med sin individuelle Beskaffenhed, og
almindelige Sætninger kan kun tilnærmelsesvis finde Anvendelse paa dem
og kan i hvert Tilfælde ikke udtømme deres Natur, ikke lade alle
Sider, alle Muligheder i dem komme til deres Ret. Men hvad vi kalder
et Menneskes Livsanskuelse er nu netop, naar den er noget værd, et
Individuelt og Person%
% --- p. 53 ---
\opage{53}ligt, bygget op paa egen Erfaring og egen Tydning, erhvervet
i Livets Kamp, gennem Arbejdet for at naa det Maal, man har sat sig.
Gennem Erfaringer, snart af Livets Storhed, Skønhed og Lykke, snart af
dets Sneverhed, Hæslighed og Lidelse, vil der uvilkaarlig forme sig et
Billede, hvis Form og hvis Farve er bestemt ved, hvad man har set,
tænkt, følt og stræbt. Den videnskabelige Erkendelse, der maatte være
tilstede, kan spille en Rolle med, men kun som et Element, der væves
ind i Billedet, ikke som altbestemmende. Hos intet Menneske, end ikke
hos den største Videnskabsmand, bestemmes Livsanskuelsen udelukkende
af Videnskab. Hos Ingen gaar Menneskeligheden helt op i
Videnskabeligheden.

Der var en Tid, da man ligesom havde Livsanskuelser paa Lager. Deres
Navne endte altid paa isme. Man havde Optimisme og Pessimisme,
Idealisme og Materialisme, Monisme og Pluralisme, Individualisme og
Socialisme o. s. v. Og man mente i Reglen ikke at behøve mange
Antydninger for at vide, hvilken af de mange Skuffer man skulde trække
ud for at anbringe en foreliggende Anskuelse paa rette Sted. Man gik
tillige ud fra, at man kendte alle mulige Retninger for menneskelig
Maade at tage Livet paa. Og dog kunde det jo være, at Livet var rigere
end som saa, og at de hidtil fundne Arter af Livsanskuelse ikke slog
til.

Der var en ung Mand i Syrakus, som gav sig Udseende af at kende
Platon's Opfattelse i dens inderste Væsen. I den Anledning udtalte
Platon, denne store, i literær Fremstilling af filosofiske Tanker
uovertrufne Mester, at hans inderste Mening ikke kunde udtrykkes i
Skrift. Kun gennem lang Syslen med Tankerne, helst under personlig
Drøftelse sammen med deres Ophavsmand, kunde fuld
% --- p. 54 ---
\opage{54}Forstaaelse naas, og der kan da pludselig dukke et Lys op,
som siden kan brænde ved egen Kraft. --- Mere eller mindre gaar det
vel saaledes, som Platon her beskriver, med al Erkendelse, men ganske
særlig gælder det om et Menneskes inderste Tanker om Livet, som det
har aabenbaret sig for ham.

18. Foruden den Modsætning, Videnskaben ved sin almene Karakter kommer
til at staa i til Livet, gør sig ogsaa et andet Modsætningsforhold mere
og mere gældende. Livsførelsen kræver Koncentration om visse bestemte
Tanker, Samling i én bestemt Retning. Men indenfor Videnskaben som
indenfor det aandelige Liv
% NOTE (p. 54): "idet-/hele" over a line end; set as one word.
idethele raader Arbejdsdeling mere og mere. Der er i den
moderne Kultur ingen intellektuel Enhedsmagt, der kan træde op overfor
Livet. Der vil da være Trang til at samle sine forskellige Erfaringer
til et individuelt Hele, selv om dette Hele ikke bygges op efter
Anvisning af Videnskab eller af Autoritet, men ved hvad Kant har kaldt
\guillemotright{}den skjulte Kunst\guillemotleft{}, der virker i os, uden at vi behøver at mærke
det. Det er en Bygning, der opføres paa den Vis, som følger af det
enkelte Individs Natur og Kaar, og som først kendes, naar det
uvilkaarlige Bygningsarbejde er skredet frem til et vist Punkt. Det er
psykologiske og historiske Forhold, ikke ren Logik, der betinger denne
Bygnings Karakter. Men Erfaring viser, at der er et vist Slægtskab
mellem Konsekvens og Karakter, som gør, at det at fastholde skrigende
Modsigelse mellem Synspunkter, der synes at fremgaa af gjorte
Erfaringer, for de fleste vil være en aandelig Lidelse, som kun fra et
asketisk Stade kan have blivende Værdi. Der vil røre sig en naturlig
Trang til at komme ud over Modsigelserne, saa at hvad der bliver
tilbage af Disharmoni, kun staar som Vidnesbyrd om Livets
% --- p. 55 ---
\setcounter{footnote}{0}%
\opage{55}Fylde og Mangfoldighed. For mange Naturer bliver der vel nok en
Brod tilbage, en vis Uro, som nu engang hører til Livets Væsen.
Giordano Bruno fremhævede i sit Skrift \guillemotright{}om de heroiske Affekter\guillemotleft{}, at
der er en indre Ild, som aldrig slukkes, saalænge der stræbes efter et
højt Maal. Heroisk Affekt opstaar netop, naar man ikke lader sig
afskrække fra at hige mod et højt Maal, fordi der er Smerte og Fare
forbunden med denne Stræben\footnote{Se nærmere om det omtalte
mærkelige Skrift af \textsc{Bruno}: \textit{Den nyere Filosofis
Historie}.\textsuperscript{3} I p. 135.}. Hvad Renaissancens mest
typiske Tænker saaledes udtalte, vil gøre sig gældende i enhver
Livsanskuelse, der ikke opbygges paa for snever Grund og ikke frygter
for at komme i Højden.

19. Foruden de i det foregaaende omtalte almindelige Forhold, der
fører til, at
% PRINTED AS IS (p. 55): "iudividuel" (turned n) for "individuel". Checked at 400 dpi.
iudividuel Livsanskuelse kommer til at spille en Rolle, ligger
der i hver Personlighed Grunde, som bevirker, at den kommer til at
lægge sig Tingene til Rette, selv om Videnskab eller Religion ikke
giver nogen Anvisning dertil. Det er Grunde, der ogsaa virker ved
Videnskabens og Religionens Tilblivelse og Udvikling, men som ligger
saa dybt i den menneskelige Natur, at de kan spores selv dér, hvor der
hverken er Tale om Videnskab eller Religion. Nogle af disse skal her
fremdrages.

I første Række maa her nævnes Ekspansionstrangen, d. v. s. den
uvilkaarlige Tilbøjelighed til at lade en Stræben, en Stemning eller en
Tankegang, som man én Gang er kommen ind paa, fortsætte og udfolde sig
videre. Det Overskud af Energi, der raades over, sættes ind paa, hvad
der saaledes har fanget Sindet, og senere Oplevelser ses og tydes i
Lyset af dem. Det er den \guillemotright{}Villiens Førstetanke\guillemotleft{}, som Dante taler
om, og som kan være saa van%
% --- p. 56 ---
\setcounter{footnote}{0}%
\opage{56}skelig at forklare, fordi dens Udspring er i det dunkle.
Istedetfor at søge og anlægge helt nye Synspunkter følges den Vej, man
nu engang er slaaet ind paa. Ogsaa i Videnskaben gaar man foreløbig ud
fra, at hvad der har vist sig at gælde i en Række af Tilfælde, gælder
for andre lignende Tilfælde. Det fører til Opstilling af en Hypotese,
som saa nærmere prøves ved at opsøge nye lignende Tilfælde eller ved
at drage Slutninger ud fra Hypotesen og se, om de stemmer med gamle og
nye Erfaringer. Det har været en Ekspansion, der (hvad enten
Fortællingen om Æblets Fald er sand eller ikke) har ført Newton til at
forsøge Faldlovens Anvendelse paa Maanens og derefter paa andre
Himmellegemers Bevægelse\footnote{Se herom \textit{Den nyere Filosofis
Historie}.\textsuperscript{3} II p. 203---205.}. Videnskaben spørger
her \guillemotright{}Hvorfor ikke?\guillemotleft{} ligesom Barnet og den Vilde. Men hvad der
stiller Videnskaben i Modsætning til Tankelivets uvilkaarlige
Udfoldelse hos os alle, er, at den formulerer Resultatet af sine
Ekspansioner saaledes at det kan være Genstand for nøjagtig Prøvelse.
Ekspansion kan ogsaa gøre sig gældende paa Følelseslivets Omraade, som
naar en dybt indgribende Glæde eller Sorg kommer til at staa som
Vidnesbyrd om, hvad man overhovedet kan vente sig af Livet, saa at
derefter enten Optimisme eller Pessimisme bliver raadende. Efter
Legenden\footnote{Oversat hos \textsc{Warren}: \textit{Buddhism in
Translation} p. 56 ff.} var det Synet af Alderdomsaffældighed, af
Sygdom og af Død, der bestemte Buddha til at trække sig tilbage fra
Verden. Disse Syn var Livserfaring nok for ham. For Pascal og
Kierkegaard var det Fortællingen om Jesus i Gethsemane, der kom til at
staa som Udtryk for, hvad Livet kunde bringe, naar man satte sig et
tilstrækkeligt højt Ideal.
% PRINTED AS IS (p. 56): "Gethsemane" ends the line with NO hyphen and
% "stemningen" opens the next (checked at 600 dpi); evidently
% "Gethsemanestemningen" with the hyphen dropped (cf. "Getsemanestemningen", p. 60).
Gethsemane stemningen blev for dem Livets egentlige, rette Grund%
% --- p. 57 ---
\setcounter{footnote}{0}%
\opage{57}stemning, som maatte raade hos alle, der vilde kalde sig efter
Jesus\footnote{Smlg. min Afhandling \textit{Pascal og Kierkegaard}.
(Tilskueren 1923).}. For andre Naturer, som Spinoza og Leibniz, kan den
Kendsgerning, at der kan opnaas Erkendelsesglæde, Glæde ved
intellektuelt Arbejde, være nok til at bestemme hele Stemningen ved
Livet. I enhver alvorlig Livsanskuelse vil der være Oplevelser, som
bliver afgørende for hele Stemningen ved Livet, og vil man forstaa en
Livsanskuelse, maa man se at finde, hvad der saaledes er blevet et
\guillemotright{}Urfænomen\guillemotleft{} for et Menneske.

Naar Homer kalder Mennesket et Væsen, der ser frem og tilbage, saa er
Forholdet dog det, at vi ser frem, før vi ser tilbage. Livet leves i
Forventning, før det leves i Erindring. Det Overskud af Kraft, hvormed
Livet begynder, tillægger fra først af Forestillinger, der fremkaldes,
en praktisk Betydning som Bebudelser, og gør dem til Motiver for
Handlen. Det er først ulykkelige, eller dog uventede, Virkninger af
Handlingerne, der én Gang for alle henviser visse Forestillinger til
Fortiden, saa at der ikke knyttes Forventninger til dem.

Expansionen fører naturlig til Expression. Den vakte Stemning kræver
Udtryk. Georg Simmel, der i sin ovenanførte Afhandling forklarer Tidens
Uro ved, at de gamle Former ikke længere er et Hjem (Gehäuse) for den
Trang, der rører sig, bruger særlig en vis Retning i vor Tids Kunst som
Exempel. Efter hans Opfattelse finder den saakaldte Futurisme sin
psykologiske Forklaring deri, at \guillemotright{}Kunstnerens indre Bevægelser
ganske, som de oplevedes, fortsættes i Værket, eller rettere som selve
Værket\guillemotleft{}. Den skal ikke støbes i nogen Form, og ingen Efterligning
skal virke med. (Der Konflikt der modernen Kultur. p. 111). Denne
Betragtning kan udvides. Hvor Livets Udfoldelse
% --- p. 58 ---
\setcounter{footnote}{0}%
\opage{58}foregaar uvilkaarlig, higer den efter Udtryk, før Oplevelserne er
drøftede gennem Eftertanke. De simpleste Domme er
Udraab\footnote{\textit{Den menneskelige Tanke}. p. 84.}, og der er
Sprogforskere, som hylder (se det sidste Kapitel i Otto Jespersens
\guillemotright{}Language\guillemotleft{}) den Hypotese, at Sprogudviklingens første Stadium er
Poesi, idet det er overvejende emotionelt betinget. Ligesom nu den
nærmere Udformning af Dommen foregaar gennem den nødvendige
Specialisering og Fremdragen af Betingelser, saaledes kommer det an
paa, om den ved Expansion ud fra en betydningsfuld Oplevelse bestemte
Anskuelse kan holde sig oppe under Livets Gang. Der er her et Punkt,
som danner en Analogi til den Bekræftelse, Newton's Expansion fik ved,
at fortsat Erfaring stemmede med, hvad Deduktion ud fra hans Hypotese
krævede. Det er naturligvis en fjern Analogi, og netop denne Fjernhed
fra det Exakte er karakteristisk for Livsanskuelsernes Kaar. Overladt
til sig selv, altsaa uden nye Expansioner ud fra lignende Oplevelser,
som den grundlæggende, vil den i Sindet vakte Bevægelse kunne vare
ved, indtil Kræfterne er udtømte. Men der kan ogsaa komme nye
Expansioner ud fra Oplevelser, af helt anden Art end de tidligere, og
der vil da kunne opstaa Delthed, Tvivl og direkte Strid i Sindet. Det
er Problemernes Tid. Det er her som overalt Diskontinuitet, der
foranlediger Problemer, medens Kontinuitet er Betingelse for Hvile og
Tryghed ved de én Gang betraadte Veje. Henri Poincaré har udtalt, at i
den videnskabelige Diskussion vil Kampen mellem Diskontinuitet og
Kontinuitet vare, saa længe Mennesket vil tænke.\footnote{\textit{Les
conceptions nouvelles de la matiére} p. 67.%
% PRINTED AS IS (p. 58 n.): „matiére“ (second reading 2026-09-30, 500 dpi).
 (I Sammelskriftet
\guillemotright{}Le matérialisme contemporain\guillemotleft{}).} Det samme gælder for
Livsanskuelsernes Vedkommende. Man
% --- p. 59 ---
\setcounter{footnote}{0}%
\opage{59}skal være \guillemotright{}tro imod sig selv og mod sin bedste Stræben\guillemotleft{},
siger en Digter. Men naar er vort Selv færdigt og modent, saa at det
blot gælder Troskab mod det, som det er blevet? Og hvilken Stræben er
den bedste? Livet kan stadig bringe os i nye Situationer, i hvilke vor
Personlighed skal bestaa en ny Prøve. Som Jaspers\footnote{\textit{Psychologie
der Weltanschauungen}.\textsuperscript{2} p. 230.} smukt har vist, er
det især i \guillemotright{}Grænsesituationer\guillemotleft{}, hvor de hidtil oplevede og fastsatte
Værdier trues med Undergang, at Personligheden skal bestaa sin Prøve.
Saadanne \guillemotright{}Grænsesituationer\guillemotleft{} er Kamp, Død, Kaos, Skyld. Til disse
Situationer er der knyttet Lidelse, og det kommer da an paa, om der er
Energi nok til at bestaa, trods den. Den udførlige Skildring, Jaspers
giver af hele dette Spørgsmaal, skyldes aabenbart Studiet af Skikkelser
som Kierkegaard og Nietzsche. (Sml. ovenfor \S~16).

Kierkegaard har opstillet en Maalestok for Aandslivets Højde, som
bestaar i den Energi, med hvilken de Modsætninger, Livet kan frembyde,
fastholdes, selv om de skulde være saa skarpe, at Livet bliver en
Lidelse. Til Grund ligger her hos Kierkegaard det fra Kant stammende
Syntesebegreb, og der opstilles da en Graderække af syntetiske
Energier, lige op til den Energi, som formaar at sejre, selv om Livets
voldsomste Modsætninger tørner mod hinanden.\footnote{Se \textit{Søren
Kierkegaard som Filosof}\textsuperscript{2} p. 119---126.} Ved
Opstillingen af denne Maalestok har Kierkegaard givet et vigtigt
Bidrag til Livsanskuelsernes Psykologi, selv om hans Anvendelse af
Maalestokken er præget af den Getsemanestemning, som nu en Gang for
ham, ligesom for Pascal, betegnede Aandslivets Højdepunkt, men som
forudsætter Afvisning af alle Opgaver, som ellers i det menneskelige
Liv kan sætte Sindet i Be%
% --- p. 60 ---
\setcounter{footnote}{0}%
\opage{60}vægelse. Hvis det er saa, at det højeste er indtraadt i Verden
paa ét bestemt Tidspunkt og paa ét bestemt Sted, og at det dér i det
afgørende Øjeblik er blevet fornægtet, endog af de nærmeste, saa er
Getsemanestemningen den eneste værdige Stemning ved Livet. Kun maa man
da ikke nøjes med en Gang om Aaret at søge den fremkaldt, men den maa
være raadende bestandig og bestemme Forholdet til alt, der ellers
kunde give Menneskelivet Værdi.

Tagen i sig selv beholder den Kierkegaardske Maalestok sin store
Betydning, selv om den ikke paa Forhaand udelukker Strid om, hvad der
staar øverst paa Værdiskalaen.\footnote{En saadan Strid førte i sin Tid
% The print letterspaces "Baggesen" and "Grundtvig" in this note.
\emph{Baggesen} og \emph{Grundtvig}. De var uenige om, hvad der var det
højeste, særlig om Erkendelsens Plads i Aandslivet, men de manglede en
Maalestok. Sml. \textit{Danske Filosofer}. p. 66---72.} Den er i
aandelig Strid tilsidst det eneste Synspunkt, man har at holde sig til.
Den Energi, med hvilken Livet føres, --- Fylden af det Indhold, der kan
sammenarbejdes, og dermed af de Modsætninger, der kan gennemleves, uden
at Livet sprænges, --- det er tilsidst, hvad det kommer an paa. Men det
vil altid være en Mangel, og tilsidst en Ulykke, naar en væsentlig Del
af Livets Indhold maa fornægtes, for at man kan sikre, hvad man kalder
det ene Fornødne. Fra et menneskeligt Standpunkt kan det ene Fornødne
netop kun være den Energi, med hvilken man tager
% PRINTED AS IS (p. 60): full stop, not comma, after "Livet" before
% lower-case "idet". Checked at 600 dpi.
Livet. idet man, alt efter sin Evne, søger en Form, i hvilken dets
Indhold kan komme til sin Ret i den enkelte Personlighed.

Der er én Ting, den antydede Maalestok
% PRINTED AS IS (p. 60): "nødvendig-/digvis" over a line end, i.e.
% "nødvendigdigvis" (the syllable "dig" set twice).
nødvendigdigvis fordrer, det er, hvad den af Carlyle indledede
Bevægelse i Retning af Personlighedsprincipet kræver: Ærlighed,
personlig Sandhed. Det gælder om ikke at tage Fejl af sin Plads paa
Livets Skala, tillægge sig selv et
% --- p. 61 ---
\setcounter{footnote}{0}%
\opage{61}Forhold til Idealer, man ikke i Aand og Sandhed hylder. Det er
naturligvis en vanskelig Opgave, der herved stilles, men den maa tages
op. Ved et enkelt Eksempel kan den belyses, naar vi tænker paa
Modsætningen mellem Tragik og Humor. Getsemanestemningen er tragisk.
Den, der paa allerførste Haand gennemlevede den, holdt Stand og gik
Lidelsens Gang til Ende. Men det forholder sig med denne Stemning, som
med al Tragik. I et tidligere Arbejde har jeg som Modsætning til den
beskrevet en Maade at tage Livet paa, som giver det store og
værdifulde det højeste Stade uden at overse Lidelserne og
Ufuldkommenhederne, og som i sin inderlige Fordybelse i det store
endog kan have et Smil til Raadighed overfor det smaa. Den
Grundstemning, jeg saaledes har søgt at karakterisere, betragter jeg
ingenlunde uden videre som den \guillemotright{}rette\guillemotleft{} eller som den \guillemotright{}bedste\guillemotleft{},
end sige som den eneste mulige, eller som den, i hvilken alle Livets
Elementer er komne til deres Ret. Jeg hævder, og det ikke blot som en
Indrømmelse, at en tragisk Livsopfattelse stedse er mulig og kan være
uundgaaelig. Det er jo muligt, at et Liv ikke kan harmoniseres uden
Fare for Indsnevring. Humoristen viger tilsidst for den tragiske Helt,
der ikke svigter sin Opgave, selv om alle Harmonier sprænges. Der er
intet i Vejen for, at Humoristen selv kan blive en tragisk Helt og
bestaa den højeste Prøve, paa hvilken Mennesker kan blive sat. Men det
er af afgørende Betydning at fastholde\footnote{Sml. \textit{Den store
Humor}. p. 104---108.}, at det Tragiske skal komme \guillemotright{}af sig selv\guillemotleft{},
ved Tingenes Gang, den Tingenes Gang, i hvilket Mennesket for at kunne
løse sin egen Opgave, har maattet gribe ind. Det er en Skæbne, som
kommer, men det er ikke Noget, som det skulde være enhvers Pligt at
opleve.

% --- p. 62 ---
\setcounter{footnote}{0}%
\opage{62}I den interessante Karakteristik, som Gorky har givet af Leo
Tolstoi, siger han, at hvad der især frastødte ham hos den store
Forfatter, var hans Attraa efter at komme til at lide for sine Ideer
for at disse derved kunde øve større Indflydelse%
% PRINTED AS IS (p. 62, note): "Reminiscenses" and "Nicolayewitch" as printed
% (checked at 400 dpi); no stop after "p. 40".
\footnote{\textit{Reminiscenses of Leo Nicolayewitch Tolstoi. By Maxim
Gorky}. London 1920. p. 40}. Et er det nemlig, at en tragisk Skæbne
kan forudses, et Andet at ønske den, naar Tingenes egen nødvendige Gang
ikke fører den med sig. Den kommer som et Jordskælv, der ikke kan
fremkaldes. En tragisk Type, kan derfor ligesaa lidt som en
humoristisk Type uden videre stilles op som Højdepunktet paa Skalaen.
Hverken Tragik eller Humor er blotte Følger af udvist psykisk Energi,
men deres Mulighed beror for en stor Del paa Forhold, der ikke staar i
vor Magt.

20. Baade overfor den Livsanskuelse, der sætter Tragik, og den, der
sætter Humor som det højeste paa Skalaen, kan Livet optræde i sin
uudtømmelige Mangfoldighed. Er det muligt at sammenfatte denne Rigdom i
én eneste Higen eller i nogensomhelst harmonisk Ordning? Der kan
stadig dukke nye Sider ved Livet op, og selv den mest anskuende
Betragtning, den største Livserfaring, kan ikke faa udtømt alt. Der vil
stedse bestaa en Brydning mellem Trang til Samling og Trang til at
udvide Horizonten. Helhedsbegrebet peger overalt, hvor det forekommer,
ud over sig selv. Det er et Grænsebegreb, som jeg har søgt at vise i
Afhandlingen om Totalitet som Kategori (Kap. 6). Det gælder ogsaa om
den Helhed, det personlige Liv stræber at blive. Og dog er det med
Rette, at den Maalestok, der anlægges paa de forskellige
Livsanskuelser, sætter Indholdets Fylde som en væsentlig Bestemmelse.
Det er let nok at sammenfatte et tarveligt Indhold.

% [text to be added: pp. 63--74]

% [text to be added: pp. 75--86]

% [text to be added: pp. 87--98]

% [text to be added: pp. 99--110]

% [text to be added: pp. 111--123]

\end{document}
